 % !TeX root = book.tex
\chapter{Web Prolog}\label{sec:web-prolog}

\begin{quotation}
\noindent Imagine a dialect of Prolog with actors and mailboxes and send and receive -- all the means necessary for powerful concurrent and distributed programming. Alternatively, think of it as a dialect of Erlang with logic variables, backtracking search and a built-in database of facts and rules -- the means for logic programming, knowledge representation and reasoning. Also, think of it as a web programming language, and as a \textit{lingua franca} for logic-based programming systems, standardised by the W3C. This is what \textbf{Web Prolog} is all about.

\flushright\textit{Web Prolog -- the elevator pitch}
\end{quotation}
\vspace{1.0cm}

\noindent %Prolog is a logic programming language. Based on formal logic, a subject dating all the way back to the antiquity and tried and tested by generations of logicians and philosophers, logic programming forms a paradigm of it own. Even in its pure form, as a set of Horn clauses, Prolog is Turing complete, and a number of extra-logical constructs has been added to ensure its usefulness as a general purpose programming language. 

\section{The essence of Prolog}

\noindent Based on formal logic, a subject dating all the way back to the antiquity and tried and tested by generations of logicians and philosophers, logic programming forms a paradigm of it own, very different from the imperative or functional programming paradigms. Prolog is generally regarded as the first logic programming language, and is arguably the most important one. Consider the following program:

\begin{verbatim}
husband(Wife, Husband) :- wife(Husband, Wife).
\end{verbatim}
\vspace{-3mm}
\begin{verbatim}
wife(socrates, xantippa).
wife(aristotle, pythias).
\end{verbatim}

\noindent We have here a \textit{rule} that translates into the following formula in predicate logic:

\begin{center}
$\forall x \forall y [wife(x,y) \rightarrow husband(y,x)]$
\end{center}

\noindent The rule in combination with the two \textit{facts} of the predicate \texttt{wife/2} (that need no translation) allow us to query the predicate \texttt{husband/2}. We may for example ask if it is true that Aristotle is the husband of Pythias, ask who is the husband of Pythias, ask who is the wife of Socrates, or enumerate the married couples one by one:

\begin{verbatim}
?- husband(pythias, aristotle).
true.
?- husband(pythias, Husband).
Husband = aristotle.
?- husband(Wife, socrates).
Wife = xantippa.
?- husband(Wife, Husband).
Wife = xantippa, Husband = socrates ;
Wife = pythias, Husband = aristotle.
?-
\end{verbatim}

% \noindent The rule in combination with the two \textit{facts} of the predicate \texttt{wife/2} (that need no translation) allow us to query the predicate \texttt{husband/2}. We may for example ask if it is true that Aristotle was the husband of Pythias:

% \begin{verbatim}
% ?- husband(pythias, aristotle).
% true.
% ?-
% \end{verbatim}

% \noindent Or we can ask who was the husband of Pythias:

% \begin{verbatim}
% ?- husband(pythias, Husband).
% Husband = aristotle.
% ?-
% \end{verbatim}

% \noindent Only one solution was found (so it appears that Pythias was not a bigamist, and nor was Aristotle):

% \begin{verbatim}
% ?- husband(Wife, aristotle).
% Wife = pythias.
% ?-
% \end{verbatim}

% \noindent Enumerate the married couples one by one!

% \begin{verbatim}
% ?- husband(Wife, Husband).
% Wife = xantippa, Husband = socrates ;
% Wife = pythias, Husband = aristotle.
% ?-
% \end{verbatim}

\noindent This simple example serves as a reminder that Prolog is not only a logic language, but also a \textit{relational} language that can be used to check if a sentence is true, to look up the value of one argument given the value of another, or to enumerate all pairs of values. When querying a binary relation, these are the \textit{modes} that exist.

%Prolog allows the use of complex terms in the arguments of clauses, here in the form of \textit{lists} in the well-known definition of \texttt{append/3} -- a predicate with the mode \texttt{append(?,\,?,\,?)}:

Because Prolog can pattern-match on structured terms in clause heads, \texttt{append/3} defines list concatenation with a base case and a recursive case that peels off the head \texttt{H} of the first list and rebuilds it on the result:

\begin{verbatim}
append([], L, L).
append([H|L1], L2, [H|L]) :- append(L1, L2, L).            
\end{verbatim}

%\noindent A query such as \texttt{?-append(Xs,\,Ys,\,Zs)} is true if \texttt{Zs} is the concatenation of the lists \texttt{Xs} and \texttt{Ys}.

\noindent All three arguments are relational -- any of them may be given or left as variables -- so \texttt{?-append(Xs,\,Ys,\,Zs)} holds precisely when \texttt{Zs} is \texttt{Xs} followed by \texttt{Ys}.

Complex terms in the arguments of clauses ensures that Prolog is a \textit{Turing complete} programming language. This is what makes it different from a language such as Datalog, which is similar in many ways, but is not Turing complete.

In order to be not only Turing complete but also a practical and efficient programming language, serious Prolog systems need to offer a lot more, and they normally do. In Section~2.1 of \textit{Fifty Years of Prolog and Beyond}, the authors provide a conceptual and minimalist definition of the important features of Prolog, and are thus able to draw a line between what can be considered a Prolog implementation and what can not. Actually, they draw \textit{two} lines, since they distinguish the \textit{essential} features of Prolog from \textit{important} (yet non-essential) features. Below, we reproduce their list of features, where 1-6 are considered essential features and 7-12 less essential. As it turns out, all the essential features except for 2) and 6) are involved when querying the binary relation between the married couples, or the ternary relation between the lists, in the examples given above.

\begin{enumerate}
\item Horn clauses with variables in the terms and arbitrarily nested function symbols as the basic knowledge representation means for both programs (a.k.a. knowledge bases) and queries;
\item the ability to manipulate predicates and clauses as terms, so that meta-predicates can be written as ordinary predicates;
\item SLD-resolution (Kowalski, 1974) based on Robinson's principle (1965) and Kowalski's procedural semantics (Kowalski, 1974) as the basic execution mechanism; 
\item unification of arbitrary terms which may contain logic variables at any position, both during SLD-resolution steps and as an explicit mechanism (e.g., via the built-in =/2);
\item the automatic depth-first exploration of the proof tree for each logic query;
\item some control mechanism aimed at letting programmers manage the aforementioned exploration;
\item negation as failure (Clark, 1978), and other logic aspects such as disjunction or implication;
\item the possibility to alter the execution context during resolution, via ad-hoc primitives;
\item an efficient way of indexing clauses in the knowledge base, for both the read-only and read-write use cases;
\item the possibility to express definite clause grammars (DCG) and parse strings using them;
\item constraint logic programming (Jaffar and Lassez, 1987) via ad-hoc predicates or specialized rules (Fruhwirth, 2009);
\item the possibility to define custom infix, prefix, or postfix operators, with arbitrary priority and associativity.
\end{enumerate}

\noindent One cannot avoid noticing that, as far as we know, no other community in support of a programming language has found itself in the position of having to determine what should be considered a \textit{real} language of this kind. As we found in Chapter~\ref{sec:intro}, the problem with Prolog is that there are so many systems around that can rightly claim to implement it but still are incompatible with each other in important ways. This is what prompted the development of a standard. For almost thirty years now, the core features of Prolog has been standardized by ISO, documented in a report known as ISO/IEC 13211-1:1995.\footnote{\url{https://www.iso.org/obp/ui/\#iso:std:iso-iec:13211:-1:ed-1:v1:en}}


\subsubsection{The logic programming landscape}
\label{sec:lp-landscape}

Prolog is the most widely known logic programming language, but it is far from the only one. Over the past five decades, researchers have extended, restricted, or reinterpreted the Horn clause core in many directions. The following survey is necessarily selective, but it covers the formalisms and systems that will reappear in later chapters when we consider how the Prolog Web can serve as an integrating substrate for the logic programming family as a whole.

Prolog's native mode of reasoning is deduction -- deriving consequences from given rules and facts -- but logic programming as a field has long explored other modes of inference as well. Abductive logic programming (ALP) reverses the direction: given an observation, it searches for hypotheses that, together with background knowledge, would explain it~\citep{Kakas1992}. Inductive logic programming (ILP) goes further still, learning new rules from examples and background theory~\citep{Muggleton1994}. Both traditions extend the reach of logic-based reasoning well beyond what a standard Prolog system provides out of the box.

Nor is deductive logic programming itself a monolith. Datalog restricts Horn clause logic to function-free programs, sacrificing Turing completeness in exchange for guaranteed termination and well-understood complexity bounds~\citep{Ceri1989,Green2013}. Well-founded semantics (WFS) addresses a different limitation, providing a three-valued account of negation in which every normal logic program has a unique minimal model assigning each ground atom the value \emph{true}, \emph{false}, or \emph{undefined}~\citep{VanGelder1991}; systems such as XSB and SWI-Prolog implement WFS via tabling, giving them stronger semantic guarantees than standard Negation As Failure. Answer Set Programming (ASP) departs more radically, searching for the stable models of a logic program rather than computing answers by resolution; a recent variant, s(CASP)~\citep{Arias2018}, combines ASP with constraint logic programming in a top-down, query-driven execution model with support for explanation generation. Probabilistic logic programming (PLP) equips logic programs with probabilistic reasoning: ProbLog~\citep{DeRaedt2007} is a Python-based toolbox that defines distributions over possible worlds and supports inference, sampling, and parameter learning, while cplint~\citep{Riguzzi2018} is a complementary SWI-Prolog library for LPADs under the distribution semantics, accessible both locally and through ``cplint on SWISH.'' Constraint logic programming (CLP) extends Prolog's unification with constraint solvers over domains such as finite integers, reals, and sets, and most major Prolog systems now ship with CLP libraries as standard. Finally, Logic Production Systems (LPS) \citep{kowalski2015reactive} combines logic programming with reactive, event-driven computation.

The common thread is that logic programming is not a single language but a family of formalisms. This breadth will become relevant in later chapters, where we show how the Prolog Web can accommodate any system in this family that can accept a goal and return answers.


%\subsubsection{The logic programming landscape}
%\label{sec:lp-landscape}
%
%Prolog is the most widely known logic programming language, but it is far from the only one. Over the past five decades, researchers have extended, restricted, or reinterpreted the Horn clause core in many directions. The following survey is necessarily selective, but it covers the formalisms and systems that will reappear in later chapters when we consider how the Prolog Web can serve as an integrating substrate for the logic programming family as a whole.
%
%\begin{itemize}
%
%\item \textbf{Datalog} is a function-free fragment of Horn clause logic that sacrifices Turing completeness in exchange for guaranteed termination and well-understood complexity bounds~\citep{Ceri1989,Green2013}. 
%
%\item \textbf{Well-founded semantics (WFS)} provides a three-valued semantics for logic programs with negation, assigning each ground atom the value \emph{true}, \emph{false}, or \emph{undefined}~\citep{VanGelder1991}. Systems such as XSB and SWI-Prolog implement WFS via tabling, giving them stronger semantic guarantees for negation than standard Prolog's Negation As Failure.
%
%\item \textbf{Answer Set Programming (ASP)} takes a model-generation approach: an ASP system searches for the stable models of a logic program rather than computing answers by resolution. A recent variant, s(CASP)~\citep{Arias2018}, combines ASP with constraint logic programming in a top-down, query-driven execution model with support for explanation generation.
%
%\item \textbf{Probabilistic logic programming} equips logic programs with probabilistic reasoning. ProbLog~\citep{DeRaedt2007} is a Python-based toolbox that defines distributions over possible worlds and supports inference, sampling, and parameter learning. cplint~\citep{Riguzzi2018} is a complementary SWI-Prolog library supporting inference and learning for LPADs under the distribution semantics, accessible both locally and through ``cplint on SWISH.''
%
%\item \textbf{Constraint logic programming (CLP)} extends Prolog's unification with constraint solvers over domains such as finite integers, reals, and sets, enabling aggressive search-space pruning. Most major Prolog systems ship with CLP libraries as standard.
%
%\item \textbf{Multi-paradigm systems.} Picat~\citep{Zhou2024} integrates logic programming with constraint solving, dynamic programming, planning, and scripting. LPS~\citep{Kowalski2016} combines logic programming with reactive, event-driven computation and is implemented in SWI-Prolog.
%
%\end{itemize}
%
%\noindent The common thread is that logic programming is not a single language but a family of formalisms. This breadth will become relevant in later chapters, where we show how the Prolog Web can accommodate any system in this family that can accept a goal and return answers.
%

\section{Web Prolog in a nutshell}\label{sec:wp-nutshell}

Web Prolog is designed as a superset of a subset of the ISO Prolog core standard. The dialect defined by ISO is well understood, and is reasonably close to most of the dialects used in Prolog textbooks. The ISO Prolog syntax specification as well as a large subset of its built-in predicates form an excellent point of departure for our attempt to create a design and a standard for Web Prolog. 

The syntax of Web Prolog is exactly as in ISO Prolog, except that three infix operators (\texttt{!/2}, \texttt{if/2} and \texttt{@/2}) that are not in the standard are defined. They can easily be added by means of \texttt{op/3}, the ISO Prolog predicate allowing a programmer to define prefix, postfix or infix operators and specify their precedences. This means that syntactically well-formed Web Prolog programs can always be \textit{read} (but not necessarily \textit{run}) by a conforming implementation of ISO/IEC 13211-1:1995. It is easy to imagine circumstances where this might come in handy. 

The capabilities of a programming language are to a large extent determined by the built-in constructs it provides. Rather than to make a long list, Figure~\ref{fig:wordcloud2} shows a ``cloud'' of built-in control constructs and predicates specified by the ISO standard. 

\begin{figure}[h]
    \centering
	\includegraphics[width=9cm]{wordcloud2}
    \caption{A cloud of ISO Prolog built-in predicates.}
    \label{fig:wordcloud2}
\end{figure}

\noindent Most of the predicates in this cloud \textit{can} and \textit{should} be supported by Web Prolog. (We will get to the exceptions a bit further ahead.) However, while most of them are necessary, they are not sufficient. Here is what else we think is needed:

\begin{itemize}

\item In addition to the built-in predicates offered by the ISO Prolog core, Web Prolog needs to support a set of \textit{standard library predicates} that conforming nodes on the Prolog Web must normally offer. They should just be there -- a programmer should not have to load them explicitly. The \textit{Prologue for Prolog} lists \texttt{member/2}, \texttt{append/3}, \texttt{length/2}, \texttt{between/3}, \texttt{select/3}, \texttt{succ/2}, \texttt{maplist/2-8}, \texttt{nth0/3}, \texttt{nth1/3}, \texttt{nth0/4}, \texttt{nth1/4}, \texttt{call\_nth/2}, \texttt{foldl/4-6} and \texttt{countall/2}.\footnote{\url{https://www.complang.tuwien.ac.at/ulrich/iso-prolog/prologue}} See also the public-domain set of libraries developed by the Prolog Commons Working Group.\footnote{\url{http://prolog-commons.org}}

\item Item number 10 in the list of essential and important features points to the possibility to express Definite Clause Grammars (DCGs) and parse strings using them. DCG notation and the associated parsing primitives are specified separately as an extension, in ISO/IEC 13211-3 (Part 3: Definite clause grammar rules).

\item As we are aiming for a profile of Prolog suitable for programming the Web we need predicates such as \texttt{http\_open/3} for accessing its resources. Also, JSON serialization/deserialization seems essential to have.

\item Last but not least, we extend our subset of ISO Prolog with carefully crafted concurrency and distribution primitives heavily inspired by Erlang. When we write ``heavily inspired,'' we really mean it. We try to stay as close to Erlang as possible, we use the same \textit{kind} of actor as Erlang, and we use the same names for our concurrency-oriented predicates that Erlang uses for the corresponding functions. This is not to say that we never deviate from the way Erlang does its thing, because as we shall see, we do when we have good reason to.

\end{itemize}

\noindent We would expect the community to be able to agree on the first three points, and the last point in the list is likely to present the only major technical challenge. It is not \textit{much} of a challenge however, since Erlang shows us where we are lacking, and where we need to go.

There are features, built-in predicates and libraries that may never make it into Web Prolog. The reason why may vary -- they may be dangerous, or they may be superfluous. Other features may be deemed not developed enough to be included at this point in time.

\begin{itemize}

\item We only use a subset of ISO Prolog as it contains primitives that do not seem to ``belong'' on the Web and which may also compromise the security of a node, such as predicates that gives a program access to the operative system.

\item Even though they are certainly relevant to the Web, there is definitely no need for predicates for building web servers, such as the ones available in SWI-Prolog's \texttt{library(http/http\_server)}. After all, a Prolog node \textit{is} a web server.

\item Even though predicates for constraint logic programming is listed as an important Prolog feature, and unless they are already defined by ISO, we do not believe time is ripe to include them in Web Prolog. As they mature, this may of course change.

\item We expect that it is possible to make only limited use of modules. After all, a node \textit{is} a module, and an actor may be seen as a module too.

\end{itemize}

\noindent We have characterized Web Prolog as a profile of Prolog suitable for programming on the Web. As we shall see, Web Prolog can be cut up into subsets that are themselves profiles, or sub-profiles if you will. Some such profiles of Web Prolog do not support predicates for database updates such as assert and retract, or predicates for reading from and writing to a terminal. Other profiles do not support \texttt{op/3}.

\medskip
\noindent FYPB argues that the ISO core standard helped establish a common kernel and improved developer confidence in portability, but also that portability problems persist and that converging on shared libraries and tools remains difficult in practice. From that perspective, Web Prolog is best seen as a portability strategy rather than a language fork: we keep the ISO syntax, constrain dangerous or non-web-centric features, and then standardise a small, web-relevant set of libraries and concurrency primitives that every node is expected to provide. This shifts the portability problem from ``porting whole systems'' to ``agreeing on profiles and their mandatory capabilities'' -- a more tractable target in an open ecosystem.

In the terminology of the paper, this also aligns with more agile convergence efforts (such as Prolog Commons) that aim to improve practical compatibility and shared infrastructure without requiring the full weight of an ISO revision process.

%\subsection{The Logic Programming Landscape}
%\label{sec:lp-landscape}
%
%Prolog is the most widely known logic programming language, but it is far from the only one. Over the past five decades, researchers have developed a rich family of formalisms that extend, restrict, or reinterpret the Horn clause core in different directions, each with its own semantic foundations, complexity properties, and practical strengths. To understand what Web Prolog aims to integrate, one must first appreciate the breadth of what logic programming has become.
%
%\begin{itemize}
%
%\item \textbf{Datalog.} Datalog is a function-free fragment of Horn clause logic: it admits only constants and variables, never compound terms. This restriction sacrifices Turing completeness in exchange for guaranteed termination and well-understood complexity bounds -- properties that matter when queries run over large, potentially untrusted datasets. Datalog has seen a sustained resurgence as a query and analysis language in domains ranging from program analysis to access control to distributed systems~\citep{Ceri1989,Green2013}. It also plays a foundational role in parts of the Semantic Web stack, where it is closely related to Description Logic Programs (DLP), a fragment designed to connect Datalog-style rule reasoning with the ontological expressiveness found in OWL~\citep{Hendler2008}.
%
%\item \textbf{Well-founded semantics.} Standard Prolog's treatment of negation through Negation As Failure (NAF) is operationally useful but semantically fragile: the order of clause evaluation can affect which negative conclusions are drawn. Well-founded semantics (WFS) provides a cleaner three-valued account in which every normal logic program has a unique minimal model, assigning each ground atom the value \emph{true}, \emph{false}, or \emph{undefined}~\citep{VanGelder1991}. Mature Prolog systems such as XSB and SWI-Prolog already implement WFS via tabling, and the combination of tabling with well-founded negation gives these systems a computational behaviour closer to the model-theoretic ideal than raw SLD resolution.
%
%\item \textbf{Answer Set Programming.} Answer Set Programming (ASP) approaches logic programming from a model-generation perspective. Rather than computing answers by resolution, an ASP system grounds a logic program and then searches for its stable models -- the answer sets -- using techniques adapted from SAT solving. ASP excels at combinatorial search, planning, and configuration problems, and mature systems such as Clingo~\citep{Gebser2019} are widely used in both research and industrial applications. Although ASP's execution model differs substantially from Prolog's top-down evaluation, the underlying language of rules, constraints, and negation places it squarely within the logic programming family. A notable recent development is s(CASP)~\citep{Arias2018}, which combines ASP with constraint logic programming in a top-down, query-driven execution model that includes support for explanation generation, bringing ASP closer to the interactive, goal-directed style familiar from Prolog.
%
%\item \textbf{Probabilistic logic programming.} A different line of extension equips logic programs with probabilistic reasoning. Systems such as ProbLog~\citep{DeRaedt2007} allow facts to hold with specified probabilities, defining a distribution over possible worlds from which the probability of any query can be computed. ProbLog~2 is a Python-based toolbox supporting inference, sampling, most probable explanations, and parameter learning, available both as a library and through online tools from the DTAI group at KU Leuven. A complementary system is cplint~\citep{Riguzzi2018}, a suite of tools for probabilistic logic programming implemented as a SWI-Prolog library. cplint supports both inference and learning for Logic Programs with Annotated Disjunctions (LPADs) and related formalisms under the distribution semantics. It is accessible both as a local Prolog pack and through ``cplint on SWISH,'' a web-based environment where probabilistic logic programs can be edited, queried, and learned from data via a browser.
%
%\item \textbf{Constraint logic programming.} Constraint logic programming (CLP) extends the unification mechanism at the heart of Prolog with constraint solvers over specific domains -- finite integers, real numbers, sets, and others. A CLP system interleaves goal reduction with constraint propagation, pruning the search space far more aggressively than pure backtracking. CLP has proved effective in scheduling, resource allocation, and combinatorial optimisation, and most major Prolog systems ship with CLP libraries as standard.
%
%\item \textbf{Multi-paradigm systems.} Some languages take a broader view of what logic programming can be. Picat~\citep{Zhou2024} is a multi-paradigm, Prolog-like language that integrates logic programming, constraint solving, dynamic programming, planning, and scripting in a single framework. Its pattern-matching rules, tabling, and constraint-based primitives make it well suited for combinatorial search, AI planning, and optimisation tasks, while its familiar syntax aims to provide a practical alternative to both traditional Prolog systems and dedicated constraint or planning languages. Another notable example is LPS~\citep{Kowalski2016}, a framework that combines logic programming with reactive and event-driven computation, implemented in SWI-Prolog and accessible through SWISH.
%
%\end{itemize}
%
%\medskip
%\noindent
%The common thread is that logic programming is not a single language but a family
%of formalisms united by a shared commitment to declarative specification and
%logical inference. This breadth will become relevant in later chapters, where we
%show how the architecture of the Prolog Web can accommodate not only classical
%Prolog but, in principle, any system in this family that can accept a goal and
%return answers.


\subsection{Web Prolog is inspired by Erlang}\label{sec:erlang-inspiration}
\vspace{-0.5cm}
\begin{quotation}
\flushright I would prefer multi-threading in Prolog to look as much as possible like Erlang.

\flushright\textit{Richard O'Keefe}\footnote{\url{https://groups.google.com/g/erlang-programming/c/1jdsnqZ4XfQ/m/ve9WfFl2YBwJ}}
\end{quotation}

\noindent Erlang is a general-purpose, concurrent, functional programming language developed by Joe Armstrong, Robert Virding and Mike Williams in 1986. Regarded as the first actor programming language to gain popularity, it started out as a proprietary language within Ericsson,\footnote{\url{https://www.ericsson.com}} but was released as open source in 1998.\footnote{A good overview of the Erlang language, its design intent and the underlying philosophy, is given in the Ph.D. thesis of Joe Armstrong \cite{Armstrong2003}.} Erlang was designed with the aim of improving the development of telephony applications, but has in recent years, thanks to its built-in support for massive concurrency, distribution and fault tolerance, been used as a very capable language for implementing the server-sides of web applications,\footnote{By companies such as WhatsApp and Klarna for example.} as well as for programming a wide range of soft real-time control problems \cite{10.5555/229883}.

Inspired by the list of essential features of Prolog given in the previous section, and if given the task of formulating a similar list defining what it means to be an Erlang-style programming language, we would likely include at least the following four essential requirements:

\begin{enumerate}
\item
  The ability to execute a large number of actor processes concurrently;
\item 
  the ability of actors to keep their states isolated so that no one else can mutate them directly;
\item
  the ability to use message-passing for asynchronous communication between actor processes, avoiding mutable shared memory and locking issues;
\item
  the ability to support network-transparent concurrent programming where actors are allowed to create other actors on remote computers and enter into seamless communication with them.
\end{enumerate}

\noindent A more succinct expression that captures the essence of Erlang is to characterize it as a language for \textit{message-passing concurrency and distribution}.

These are four essential features that if added to the features of ISO Prolog will provide us with the most crucial parts of a foundation for Web Prolog, and indeed for the whole Prolog Trinity ecosystem. Of course, for this to work we must make sure that the two sets of features are compatible. Our current understanding is that they are, and we intend to demonstrate this in the book.

Why did we choose Erlang as a source of inspiration, rather than any alternative? After all, there are other actor programming languages -- Akka, Pony and E, to name a few. The main reason for the choice of Erlang is that it is arguably the most \textit{mature} concurrency-oriented language in existence, with a bigger community than the alternatives, and a community that includes a fair number of industrial users. Many millions of lines of source code have been written in Erlang, many books teaching the language have been authored, and university courses teaching concurrent and distributed programming often uses Erlang. In our opinion, since we are aiming for a standard it makes a lot of sense to borrow the required concurrency-oriented features from a language as mature and battle-tested  as Erlang.

Another reason for our choice is that Erlang and Prolog are in many ways remarkably similar, both when it comes to syntax, and the reliance on dynamic typing, assign-once variables, pattern matching and recursion. The similarities can be explained by the fact that historically, Erlang evolved from Prolog and the first version of Erlang was actually implemented as an interpreter in Prolog \cite{useofprolog:article}. %With time, it evolved into a simple functional language.
The following listings of Erlang code (to the left) and Web Prolog code (to the right) show some of the similarities as well as some of the differences between the two languages:

\begin{figure}[h]
\begin{minipage}{0.485\linewidth}
\begin{tcolorbox}[
  colback=white,
  colframe=black,
  arc=0pt,
  outer arc=0pt,
  boxrule=0.4pt,
  left=2mm,right=2mm,top=1mm,bottom=1mm,
  enhanced
]
%\small
\begin{verbatim}
append([], L) -> L ;
append([H|L1], L2) -> 
    [H|append(L1, L2)].                
\end{verbatim}
\end{tcolorbox}
\end{minipage}
\hfill
\begin{minipage}{0.485\linewidth}
\begin{tcolorbox}[
  colback=white,
  colframe=black,
  arc=0pt,
  outer arc=0pt,
  boxrule=0.4pt,
  left=2mm,right=2mm,top=1mm,bottom=1mm,
  enhanced
]
%\small
\begin{verbatim}
append([], L, L).           
append([H|L1], L2, [H|L]) :-
    append(L1, L2, L). 
\end{verbatim}
\end{tcolorbox}
\end{minipage}
\end{figure}


% \begin{verbatim}
% append([], L, L).                append([], L) -> L ;
% append([H|L1], L2, [H|L]) :-     append([H|L1], L2) -> 
%     append(L1, L2, L).               [H|append(L1, L2)].
% \end{verbatim}

\noindent Note the use of capital letters for variables, the familiar notation for lists, and the reliance on pattern matching and recursion. A major difference is that Erlang is a functional programming language, whereas Web Prolog is relational. Since a function is just a special case of a relation this difference must not be exaggerated, but it does show in the way function calls may be nested, something that cannot be done in Prolog (or Web Prolog) since arguments are used to represent outputs as well as inputs. What \textit{can} be done in Prolog, however, but not in Erlang, is to call \texttt{append/3} in more than one \textit{mode} -- not only to append one list to another, but also, for example, to nondeterministically split a list up into two parts. Furthermore, thanks to logic variables, \texttt{append/3} is tail-recursive in Prolog, but not in Erlang's \texttt{append/2} (as it is written here).

Despite such differences, as long as we do deterministic computation only, and no search is involved, logic programming and functional programming are fairly similar in the way they work, and methods used to achieve success with one often transpose to the other. Again, features such as pattern matching, assign-once variables and recursion play important roles in both languages.

During its evolution, Erlang gradually shed several features that Prolog had -- and still has -- while at the same time developing exceptionally strong support for concurrent and distributed programming, an area where Prolog historically lacked, and still lacks, comparable built-in capability. Key here is the concept of an actor, and the ways actors can be scripted. Below, the echo server written in Web Prolog is placed side-by-side with an implementation in Erlang -- a predicate in Web Prolog and a function in Erlang -- both making a recursive call that implements a loop: 

\begin{figure}[h]
\begin{minipage}{0.485\linewidth}
\begin{tcolorbox}[
  colback=white,
  colframe=black,
  arc=0pt,
  outer arc=0pt,
  boxrule=0.4pt,
  left=2mm,right=2mm,top=1mm,bottom=1mm,
  enhanced
]
%\small
\begin{verbatim}
echo_actor() ->           
    receive                
        {echo, From, Msg} ->
            From ! {echo,Msg},
            echo_actor()  
    end. 
\end{verbatim}
\end{tcolorbox}
\end{minipage}
\hfill
\begin{minipage}{0.485\linewidth}
\begin{tcolorbox}[
  colback=white,
  colframe=black,
  arc=0pt,
  outer arc=0pt,
  boxrule=0.4pt,
  left=2mm,right=2mm,top=1mm,bottom=1mm,
  enhanced
]
%\small
\begin{verbatim}
echo_actor  :-           
    receive({               
        echo(From, Msg) ->  
            From ! echo(Msg),
            echo_actor     
    }).                          
\end{verbatim}
\end{tcolorbox}
\end{minipage}
\end{figure}

% \begin{verbatim}
% echo_actor  :-                echo_actor() ->           
%     receive({                      receive                
%         echo(From, Msg) ->             {echo, From, Msg} ->
%             From ! echo(Msg),              From ! {echo,Msg},     
%             echo_actor                     echo_actor()  
%     }).                            end.                                 
% \end{verbatim}

\noindent The Web Prolog code on the left demonstrates the use of two constructs foreign to traditional Prolog, implementing the sending and receiving of messages in the style of Erlang. The programs have a similar look, and this is intentional. We have \textit{strived} to make Web Prolog look as similar as possible to Erlang within the constraints imposed by the syntax and semantics of ISO Prolog.\footnote{Note that \texttt{!/2} here is the Erlang-style \textit{send} operator, rather than the Prolog cut (\texttt{!/0}).} 

A function call such as in \texttt{Pid\,=\,spawn(fun()\,->\,echo\_actor()\,end)} (or, but only for zero-argument functions, \texttt{spawn(fun echo\_actor/0)}) can be made in order to spawn our echo server in Erlang, while doing it in Web Prolog would use something like \texttt{spawn(echo\_actor,\,Pid)}. These calls look somewhat similar too, but while Erlang is a higher-order language in which the spawn function takes an anonymous function as its argument, Prolog (or Web Prolog) is not a higher-order language in this sense. In Web Prolog, \texttt{spawn/2} is a \textit{meta predicate} which expects a callable term to be passed in the first argument, a term that when called will execute a procedure that determines how the actor will behave.

\medskip
\noindent An important observation in FYPB is that when portability problems are not solved but merely worked around, the community tends to fragment into smaller, incompatible subgroups, making shared tooling and shared libraries harder to sustain. This matters here because concurrency and distribution are precisely the kinds of features that, if left to ad-hoc, system-specific designs, will amplify dialect drift. One motivation for staying close to Erlang -- even down to predicate names and a small set of primitives -- is therefore to reduce the degrees of freedom. If we are going to add distributed concurrency to Prolog, it is better to do so in a way that maximises the chance of convergent implementations.


\subsection{Web Prolog is a multi-paradigm programming language}\label{sec:p-language}

\begin{quote}
\flushright Prolog is different, but not that different.
\vspace{-0.2cm}
\flushright \textit{Richard O'Keefe}
\end{quote}
\vspace{0.2cm}

\noindent A common way to categorize programming languages is in terms of the programming paradigm(s) they support. In the very nicely organized and very well argued taxonomy of programming paradigms by Peter van Roy in \cite{van2009programming}, depicted in Figure~\ref{fig:paradigms},\footnote{The diagram is available at https://www.info.ucl.ac.be/\textasciitilde{}pvr/paradigmsDIAGRAMeng108.pdf} it is indeed possible to pinpoint exactly where Web Prolog fits in. We have enclosed the relevant square boxes in boxes with rounded corners. 

\begin{figure}[h]
    \centering
	\includegraphics[width=\textwidth]{paradigms.png}
    \caption{A taxonomy of programming paradigms due to Peter van Roy. (Image source and license to be confirmed.)}
    \label{fig:paradigms}
\end{figure}

This suggests that we can regard Web Prolog as a language that supports three programming models: \textit{relational logic programming} (like in Prolog), \textit{imperative programming} (also like in Prolog), and (like in Erlang) \textit{message-passing concurrent programming}. (A note on terminology is in order. In this book we prefer the term \textit{actor-based programming}. Joe Armstrong used to refer to it as \textit{concurrency-oriented programming}, or COP.).  

With access to predicates for asserting and retracting clauses in the dynamic database, traditional Prolog has always, albeit somewhat reluctantly,  been able to support imperative programming (and later additions such as \texttt{setarg/3} has of course amplified this). Also, the built-in backtracking \textit{search} that is so useful in Prolog has little to do with logic. Thus Prolog has never really been a single-paradigm programming language, and this is reflected in van Roy's taxonomy. In \cite{van2009programming} he suggests that this is not a bad thing:

\begin{quote}
A good language for large programs must support several paradigms. One approach that works surprisingly well is the \textit{dual-paradigm language}: a language that supports one paradigm for programming in the small and another for programming in the large. 
\end{quote}

\noindent With the addition of support for the actor programming model, Web Prolog becomes a much clearer as well as (in our view) a more interesting case of a multi-paradigm programming language. In particular, it can be argued that the support for programming in the large has been greatly improved.

\subsubsection{A sip of Elixir and the rise of Phoenix}\label{sec:a-sip-of-elixir}

\begin{quotation}
\noindent The rise of popularity of the Internet and the need for non-interrupted availability of services has extended the class of problems that Erlang can solve.
\vspace{-0.3cm}
\flushright\textit{Joe Armstrong}
\vspace{0.3cm}
\end{quotation}

\noindent Erlang has successfully been used on the server sides of many web applications. However, this is also an area of application which in recent years has been taken over by the Elixir programming language.\footnote{\url{https://elixir-lang.org/}} A program written in Elixir compiles into the Erlang VM and therefore belongs in the same paradigm and shares the same abstractions for building concurrent, distributed and fault-tolerant applications. The main difference concerns syntax.

The syntax of Elixir is inspired by the Ruby programming language. Here is how the programs in Section~\ref{sec:erlang} looks like when written in Elixir:

\begin{verbatim}
# length

def length([]), do: 0
def length([_h|t]), do: 1 + length(t)

# naive reverse

def reverse([]), do: []
def reverse([h|t]), do: reverse(t) ++ [h]

# spawning example

self = self()
spawn fn -> send(self, length([:a,:b,:c])) end
receive do
    m -> IO.puts(m)
end
\end{verbatim}

\noindent The Phoenix web framework\footnote{\url{https://phoenixframework.org}} is an add-on to Elixir which is inspired by Ruby on Rails,\footnote{\url{https://en.wikipedia.org/wiki/Ruby_on_Rails}} a server-side MVC-based web application framework written in Ruby that was very popular only a decade ago and which can be said to have set the standard for such frameworks. In addition, and more importantly, Phoenix takes advantage of the Erlang VM's ability to handle millions of processes and connections in order to offer bidirectional real time capabilities with channels between JavaScript on the client and Elixir on any of the servers in a cluster. Each single visitor to a website can have its own process on the server and its own real time connection. This opens up possibilities that are not present in traditional web frameworks.

Note that it is true also for the Prolog Web that each client can have its own process on the server (i.e. on the node) and its own real time connection. It will not scale as well as with Elixir+Phoenix, but this might improve in the future.

Finally, keep in mind that the Prolog Web is not a framework, but an extension of the Web. Elixir+Phoenix is not webized in the same strong sense as Web Prolog and the Prolog Web, and its creators do not aim for it to become an official web standard.\\


\subsubsection{Is Web Prolog an object-oriented language?}\label{sec:OOP}

Object-oriented features are often presented as a practical route to programming in the large. In \textit{Fifty Years of Prolog and Beyond}, the lack of object orientation is identified as one of Prolog's weaknesses, and the integration of logic programming with object-oriented features is described as an ``elusive goal.'' Nevertheless, the authors present a number of what they call ``effective'' proposals for object-oriented extensions to Prolog, developed by various researchers and system builders. Among these, Paolo Moura's \textit{Logtalk} stands out as particularly influential. 

Interestingly, Alan Kay -- often credited with inventing object-oriented programming (OOP) -- seems to regard message passing as the fundamental concept in object-oriented programming, rather than the more visible constructs of classes and objects. He even considers \textit{Erlang} itself to be an object-oriented language:\footnote{\url{https://computinged.wordpress.com/2010/09/11/moti-asks-objects-never-well-hardly-ever/}}

\begin{quotation}
Significant parts of Erlang are more like a real OOP language than the current Smalltalk, and certainly the C based languages that have been painted with ``OOP paint.''
\end{quotation}

\noindent What is evident is that \textit{if} Erlang qualifies as an object-oriented programming language, then Web Prolog too must be considered one. And \textit{if} so, this means that Web Prolog should probably be added to the list of object-oriented extensions of Prolog referenced by the authors of FYPB. But those are \textit{big} ifs, and in a blog post Armstrong writes:\footnote{\url{https://harmful.cat-v.org/software/OO\_programming/why\_oo\_sucks}}

\begin{quotation}
As Erlang became popular we were often asked ``Is Erlang OO'' -- well, of course the true answer was ``No of course not'' -- but we didn't to \emph{[sic]} say this out loud -- so we invented a serious of ingenious ways of answering the question that were designed to give the impression that Erlang was (sort of) OO (If you waved your hands a lot) but not really (If you listened to what we actually said, and read the small print carefully).
\end{quotation}

\noindent So we still prefer to refer to Web Prolog processes as ``actors'' and ``agents,'' rather than ``objects,'' and to frame the paradigm as ``actor-programming'' or ``agent-oriented programming'' rather than ``object-oriented programming.'' Programmers looking for an OO-system that is class-based and supports inheritance would likely prefer Moura's Logtalk, while the rest of us may be happy with the features provided by Web Prolog.

\subsubsection{Any unexpected interactions between paradigms?}

It is probably wise to entertain a suspicion of unexpected interactions between language features and possible impedance mismatches between the two paradigms -- between Prolog's relational, nondeterministic programming model based on logic and Erlang's functional and message passing model. How well do the Erlang-style constructs mix with Prolog -- with backtracking for example, or with the features for imperative programming? What do we get if we combine them? A kludge, or something quite beautiful? (This, as so many other things, might be in the eye of the beholder, but we know what \textit{our} eyes tell us.)

In theory, we should be on the safe side. Sequential Erlang is basically Erlang with its data types, one-way pattern matching, functions and control structures, but without spawn, send, receive, and other constructs used for concurrent programming. The idea behind Web Prolog can thus be described as an attempt to ``plug out'' the sequential part from Erlang and ``plug in'' sequential Prolog instead. There seems to be no principled reasons why we would not be able to replace the sequential functional language with a sequential relational language, e.g. a logic programming language such as Prolog. Indeed, Joe \citet{armstrong2003concurrency} writes:

\begin{quotation}
\noindent Erlang is a concurrent programming language with a functional core. By this we mean that the most important property of the language is that it is concurrent and that secondly, the sequential part of the language is a functional programming language.

The sequential subset of the language expresses what happens from the point it time where a process receives a message to the point in time when it emits a message. From the point of view of an external observer two systems are indistinguishable if they obey the principle of observational equivalence. \textit{From this point of view, it does not matter what family of programming language is used to perform sequential computation.} (Our emphasis.)
\end{quotation}

\noindent Alternatively -- and this has been our choice -- the approach can be described as an attempt to \textit{extend} Prolog with constructs such as spawn, send and receive. The idea is to keep everything that core Prolog has to offer, and extend it with a number of those primitives that make Erlang such a great language for programming message-passing concurrency. The choice between extending Prolog with Erlang-style constructs and extending Erlang with Prolog-style constructs is easy to make, and a lot has to do with syntax.
Provided we can accept using a syntax which is relational rather than functional, precluding the nesting of function calls, it may be clear already at this point that the surface syntax of Prolog can easily be adapted to express the needed Erlang-style primitives. It is arguably a lot harder to express Prolog rules and other constructs using the syntax of Erlang.

\subsection{Web Prolog as a scripting language for the Web}\label{sec:web-prolog-as-a-scripting-language}

\begin{quotation}
\noindent Scripting languages are a lot like obscenity. I can't define it, but I'll know it when I see it.
\flushright\textit{Larry Wall}
\vspace{0.3cm}
\end{quotation}

\noindent Over the years, the Web has become a key delivery platform for a variety of sophisticated interactive applications in just about any conceivable domain. As a consequence, JavaScript, more or less the only game in town for programming the front-end of a web application, has become among the most commonly used programming languages on Earth. 
%It appears that even back-end developers are more likely to use it than any other language.\footnote{\url{http://stackoverflow.com/research/developer-survey-2016\#most-popular-technologies-per-occupation}} 
Indeed, JavaScript must be regarded as \textit{the} web programming language of our times. 

JavaScript started out as a very small, highly domain-specific scripting language that was limited to running within a web browser to dynamically modify the web page being shown, but has over time evolved into a widely portable general-purpose programming language. Modern JavaScript is a high-level, dynamically typed, interpreted multi-paradigm programming language with several essential features that make it a versatile tool for web development. JavaScript\textquotesingle s syntax supports various programming paradigms, including imperative, functional, and object-oriented programming styles. With features such as callbacks, promises, and async/await, JavaScript supports asynchronous programming, enabling non-blocking operations, especially useful in web applications. JavaScript\textquotesingle s native format for data interchange, JSON, is lightweight and widely used for data transmission. JavaScript runs on almost all modern web browsers and platforms, making it universally applicable for web development. It easily interfaces with numerous web APIs for tasks like accessing the Document Object Model (DOM), making HTTP requests, and handling events. This makes it an excellent tool for connecting disparate systems in web applications.

With the advent of Node.js,\footnote{\url{https://en.wikipedia.org/wiki/Node.js}} JavaScript extended its reach to server-side development. This allowed for a unified language experience across both client and server sides, simplifying the development process by allowing the same language to be used throughout the entire stack. Not having to learn more than one language in order to become a so called ``full-stack developer,'' and not having to switch languages when changing from working on the server-side to working on the client-side are important advantages. Also, Node.js is frequently used to build and connect microservices. Its non-blocking I/O model and event-driven architecture make it suitable for building scalable and efficient network applications.

It is clear that as a relational logic programming language with the essential and important features listed above, Prolog is very different from JavaScript in terms of paradigms supported. The addition of Erlang-style concurrency and asynchronous communication does not really narrow the gap, as those are features that work differently in Erlang. What matters here is the \textit{role} that JavaScript plays on the Web. We want Web Prolog to play a similar role in the Prolog Trinity ecosystem as JavaScript does in the JavaScript ecosystem, just as clear-cut.

%---
%
%Paradigms alone hardly capture all there is to a practical programming language. Languages may also be related by having a special purpose in common, or by constraints imposed upon them by a particular area of application. They are related by family resemblances rather than anything else, and connected by a series of overlapping similarities, where no one feature is common to all languages. 
%
%---

The connection between scripting languages and the Web has been noted before. For example, in his article ``Scripting: higher level programming for the 21st Century,'' John \citet{660187} writes (about the Internet rather than the Web, but that makes little difference):

\begin{quote}
The growth of the Internet has also popularized scripting languages. The Internet is nothing more than a gluing tool. It does not create any new computations or data; it simply makes a huge number of existing things easily accessible. The ideal language for most Internet programming tasks is one that makes it possible for all the connected components to work together; that is, a scripting language.
\end{quote}

\noindent So what about the suitability of Web Prolog for programming the Web? Can Prolog challenge JavaScript in this space, at least for some applications, such as web-based AI? For this to be possible, we must ensure that Web Prolog is

\begin{itemize}
\item viable as a scripting language,
\item suitable for programming the Web, 
\item possible to implement in browsers, and
\item secure. 
\end{itemize}

\subsubsection{Web Prolog as a scripting language}\label{sec:web-prolog-as-a-scripting-language-1}

\noindent As regards the viability of Web Prolog as a scripting language we note that while hardly a definition, \citet{660187} characterizes scripting languages as follows:

\renewcommand{\theenumi}{\arabic{enumi}}
\begin{enumerate}

\item They are suitable for ``programming in the large,''

\item they are designed for ``glueing'' applications together,

\item they tend to be typeless, thus making it easier to connect components,

\item they are usually interpreted, thus providing rapid turnaround during development by eliminating compile times,

\item they are higher level than a system programming language in the sense that a single statement does more work on average, and

\item they allow rapid prototyping.

\end{enumerate}
\renewcommand{\theenumi}{\arabic{enumi}}

\noindent JavaScript does indeed fit this characterization, but given the development described above, it should now probably be described as general-purpose languages which also happens to be suitable for scripting. By contrast, languages such as C++ and Java are not well suited to that purpose..

What about Prolog? Does it satisfy Ousterhout's characterization? The points 3-6 in his list are probably true of the Prolog programming language in general. (With the possible exception of 4 -- some systems compile, but do it very quickly.) Whether or not the points 1 and 2 applies varies between systems. For example, the people behind the commercial SICStus Prolog appears to regard Prolog more as a language for writing problem-solving modules to be embedded into other code, written in Java, C++, Python or what have you.\footnote{Mats Carlsson, main developer of SICStus Prolog, personal communication} It is only a guess, but we suspect that the people behind SICStus Prolog would not be willing to characterise it as a scripting language, although that may simply be a business decision more than anything else. In contrast, acting as a glue language and a language for programming in the large is the ambition of SWI-Prolog, witness the following quote from the online manual:\footnote{\url{http://www.swi-prolog.org/pldoc/man?section=swiprolog}}

\begin{quote}
SWI-Prolog positions itself primarily as a Prolog environment for `programming in the large' and use cases where it plays a central role in an application, i.e., where it acts as `glue' between components. 
\end{quote}

\noindent Interestingly, the people behind Erlang also think of it as a glue language: 

\begin{quote}
We use Erlang as the glue to handle all orchestration, and then we use Python, C, Julia, ... It is actually a language \textit{intended} to act as a hub towards other languages. The interfaces could be protocols, could be RESTful APIs, or other programming languages. It's ideal for that.\footnote{See \url{https://www.youtube.com/watch?v=K8nxTSPHZhs5}, 8:58 into the discussion.}
\end{quote}

\noindent Erlang-style, network-transparent message-passing concurrency seems ideal for programming in the large on the Web, and for orchestrating both local and remote processes. 
%But recall van Roy's assertion that a good language for large programs must support several paradigms: one for programming in the small and another for programming in the large. By adding Erlang-style message-passing concurrency to Prolog's logic-programming core, supplemented by its imperative features, we end up with three paradigms. This cannot be a bad thing. 
Erlang glue is likely a different kind of glue, with different properties, than Prolog glue. There is hope that combining Prolog and Erlang might yield an even stronger and more flexible glue of the required kind -- a \textit{superglue}, if you will.

Despite the good fit with Ousterhout's characterization, neither Prolog nor Erlang are usually \textit{advertised} as scripting languages, and they were hardly \textit{designed} for the purpose of glueing applications together (or at least Prolog was not). They are both full-blown very flexible general-purpose programming languages with ``batteries included.'' Like JavaScript, they should probably be regarded as general-purpose languages which also \textit{happens} to be suitable for scripting. %Again, not all general-purpose languages have the properties required for that but Prolog arguably does.

Web Prolog is a very \textit{general} special-purpose scripting language. Similar to most other scripting languages, Web Prolog can be used for purposes for which ``scripting'' is not really the right word. Web Prolog comes with a focus on web \textit{logic} programming as well as web \textit{agent} programming, and in client-side browsers as well as server-side. Using Web Prolog, the owner of a node is for example able to build knowledge bases consisting of many millions of clauses -- facts as well as rules, and/or web agents that can make use of such knowledge bases.

\subsubsection{Is Web Prolog suitable for programming the Web?}\label{sec:wp-suitable-for-web}

To demonstrate that Web Prolog is suitable for programming the Web is what this book is about. But we are going further than that, as we are also aiming for a novel extension of the Web -- the Prolog Web and even an Agentic Web -- and a complete ecosystem, for which Web Prolog will serve as the default language.

\subsubsection{Web Prolog running in browsers}\label{sec:web-prolog-running-in-browsers}

Regarding our third point, for Web Prolog to become a viable web programming complement to JavaScript it is vital that it can be implemented in browsers too, since this comes with a number of advantages:

\begin{enumerate}

\item When network connection is slow, it is best to perform the majority of computations in the browser.

\item This is where, in most scenarios, the \textit{state} of an interaction between a user and an application should preferably be represented.

\item Client-side computation reduces the demands placed on nodes.

\item This is also where we find browser APIs that lets a web developer manipulate the DOM, store data locally, and add features such as spoken interaction, video or graphics to an application.

\end{enumerate}

\noindent For Web Prolog in the browser to play the same role as JavaScript in the browser currently does, it must allow Web Prolog source code to be loaded from any server on the Web by means of the HTML \texttt{<script>} element, inline or with a link. 

We do not, however, go all the way in our attempt to replicate Javascript's abilities as a scripting language. One thing that we do not at this time aim for is to make Web Prolog into a language with primitives for scripting the Document Object Model (DOM) in web browsers. With time it may well be interesting to develop such capabilities, but we believe it is simply too early to try to standardize them. In fact, we suspect that JavaScript will reign supreme in that role, and that if a browser is running Web Prolog locally, it will be as a \textit{web worker},\footnote{\url{https://en.wikipedia.org/wiki/Web\_worker}} with which the main JavaScript process is talking.

\subsubsection{Web Prolog and security}\label{sec:web-prolog-and-security}

As for security, similar to JavaScript, Web Prolog is a \textit{sandboxed language}, open to the execution of untested or untrusted source code, possibly from unverified or untrusted clients without risking harm to the host machine or operating system. Therefore, Web Prolog does not include predicates for file I/O, socket programming or persistent storage, but must rely on the host environment in which it is embedded for such features.

Web Prolog does indeed share some of the properties that made JavaScript succeed on the Web. A visit to a web application server starts a JavaScript process on the client, running code that has been downloaded from the application's host to the user's client. Such a process must be allowed to run there (i.e. the owner of the client must allow it), and when it is (and most users allow it by default), it must execute in a way that does not harm the client. When a Web Prolog process is  run on a node it must be allowed by its owner to do so, and it must run without harming the node. So one thing they share, Web Prolog and JavaScript, is the ability to run untested and untrusted source code, authored by unverified and untrusted programmers, in a sandbox on someone else's computer.  

In summary, while JavaScript\textquotesingle s initial role was confined to client-side scripting in web browsers, its capabilities have expanded
significantly. Today, it serves as a versatile glue language that
connects various components in both client and server environments, as
well as in the broader web development ecosystem. When designing Web Prolog and the rest of the Prolog Trinity ecosystem, we must of course try to learn from such a popular and versatile tool for web development, with a focus on the good parts in particular.


\section{Erlang-style message-passing concurrency in Web Prolog}

\begin{quotation}
\noindent In a nutshell, if you were an actor in Erlang's world, you would be a lonely person, sitting in a dark room with no window, waiting by your mailbox to get a message. Once you get a message, you react to it in a specific way: you pay the bills when receiving them, you respond to Birthday cards with a ``Thank you" letter and you ignore the letters you can't understand.

\flushright\textit{Fred H\'ebert}
\vspace{0.1cm}
\end{quotation}
 
% By design, Web Prolog features the \textit{same} kind of actors as those supported by the Erlang programming language. 
% Inspired by Erlang's notion of an actor, the \textit{Prolog actor} is the fundamental unit of computation in the Prolog Trinity ecosystem. The diagram in Figure~\ref{fig:actor-2} shows three actors, one of which has just been spawned, and one that has been opened up so that we can have a look inside (where on this level of abstraction they all look the same).

% \begin{figure}[h]
%     \centering
% 	\includegraphics[width=5.6cm]{an-actor-without-db}
%     \caption{The anatomy of a Prolog actor.}
%     \label{fig:actor-2}
% \end{figure}

% \noindent A Prolog actor is a \textit{process}, capable of communicating with the world around it through messaging. An actor has a unique address -- a \textit{process identifier}, or \textit{pid} for short. If we have the pid of an actor, we can message it. %Actors are effectively first-class citizens in Web Prolog because their pids can be freely assigned to variables, passed around, returned in solutions to goals, and stored in data structures, enabling powerful concurrency patterns. 
% Since pids can themselves be part of a message, it allows other actor processes to communicate back. Furthermore, an actor has a \textit{mailbox} that stores incoming messages. There is a \textit{receive} operation that allows the actor to select and process messages from the mailbox, and a \textit{send} operation that can be used to send messages to any other actor. Finally, there is a \textit{spawn} operation that allows an actor to create other actors -- child actors as it were. 

\noindent Corresponding to the send, receive and spawn operations there are built-in Web Prolog predicates such as \texttt{!/2}, \texttt{receive/1-2} and \texttt{spawn/1-3}. In this section, the use of these and other predicates will be demonstrated by examples. We choose to work with very simple examples, many of which are borrowed from tutorials and text books teaching Erlang to beginners. The major motivation for having it in this way is to try to satisfy two kinds of readers who might want to approach the material in two different ways. Readers who have a year or two of Prolog programming under their belt, but feel that time is ripe to have a look at concurrent programming, may want to look at the examples very carefully and perhaps work through them themselves using the Trinity Prolog implementation.

Readers who are very experienced Prolog programmers and teachers may also want to ask themselves if this a good way to teach students of Prolog some concurrent programming techniques, or if something else might work better. Their role is as evaluators of potential teaching material rather than learners. We know what we believe; we think Web Prolog might be suitable for teaching both Prolog-style logic programming and Erlang-style actor programming in the same system, and preferably in a web-based playground -- a great way to create as little hassle for a teacher as possible.

Then, of course, we expect the latter kind of reader to evaluate our proposal and to determine if Erlang-style concurrent programming makes sense in the context of Prolog. (We know what we think -- it does.)


\subsection{Programmer talking to actor, actor talking to itself}

If we need to hold a Prolog-style conversation with an actor, we may want to talk to a Prolog \textit{shell}. By using a terminal attached to a shell, we can interact with it the way we normally interact with Prolog, by querying it, updating its dynamic database, and so on. For example, here is how we instruct the shell to split a list up into a prefix and a postfix using \texttt{append/3}: 

\begin{verbatim}
?- append(Xs, Ys, [a,b,c]).
Xs = [], Ys = [a, b, c] ;
Xs = [a], Ys = [b, c] ;
Xs = [a, b], Ys = [c] ;
Xs = [a, b, c], Ys = [] ;
false.
?-
\end{verbatim}

\noindent This book, as we have warned, is not a beginner's textbook, so how Prolog is able to come up with solutions to such queries is not something we will cover. Our focus is rather on the concurrency-oriented, message-passing features of Web Prolog, so we begin instead by introducing a couple of simple messaging primitives. But first, using \texttt{self/1}, we need to determine the identity of the shell we will be talking to:

\begin{verbatim}
?- self(Self).
Self = 85234512.
?-        
\end{verbatim}

\noindent What we got back here is a \textit{pid}, an unforgable and locally (but not necessarily globally) unique identifier. Our proposal here is to use random integers of a size that makes it impossible in practice for anyone to \textit{guess} the pid of an actor. %\footnote{We use integers with 8 digits in the book, but 10 digits might be better.}

As stated above, if we know the pid of an actor process we can send it a message. Just as any other kind of actor, the shell is equipped with a mailbox, and this is where the message will end up. The syntax and semantics of the send operator is easy to explain. With a call such as \texttt{Actor\,!\,Message} the term \texttt{Message} is sent to the mailbox of the process identified as \texttt{Actor}, either by means of a pid or (as we shall see) using a registred \textit{name}. For example, using \texttt{!/2} with the pid of our shell we can instruct it to send \textit{itself} a message:

\begin{verbatim}
?- 85234512 ! hello.
true.
?-
\end{verbatim}

\noindent Now we can use \texttt{receive/1} to make sure that the message we sent was received by the shell and is present in its mailbox. A single receive clause with a variable in the head and \texttt{true} in the body does the job:

\begin{verbatim}
?- receive({Message -> true}).
Message = hello.
?-
\end{verbatim}

\noindent This is one of the simplest use of the receive operation we can think of, but \texttt{receive/1} really is more complex than this, so expect more to come further ahead in this chapter. One more thing about the above call to \texttt{receive/1} is worth a mention already at this point though: had the mailbox been empty, \texttt{receive/1} would have \textit{blocked} the execution of the process running as a shell until a message shows up.

There is an even simpler call to \texttt{receive/1} that one can make, namely this:

\begin{verbatim}
?- receive({}).
\end{verbatim}

\noindent This call does not wait for a any specific message, which just means that it will block forever, or until the the process that runs it is killed. Control-C will work and send the user back to the \texttt{?-} prompt.

\subsection{Two utility tools for programming in the shell}

\noindent Copying and pasting pids is not very convenient. Instead, we can use a small shell feature (borrowed from SWI-Prolog): if a variable name is prefixed with a dollar sign, the shell substitutes it with the variable's most recent binding before running the query. For example:

\begin{verbatim}
?- $Self ! hello.
true.
?-
\end{verbatim}

\noindent Here, the shell kept track of the most recent binding of \texttt{Self} to a value (\texttt{85234512} in this case) and substituted the term \texttt{\$Self} with that value before the query was run.

\medskip

\noindent During interactive development, it is also important to inspect the shell's mailbox. Using \texttt{receive/1} is often awkward here: if the mailbox happens to be empty, \texttt{receive/1} will block, effectively freezing the session until the call is interrupted (e.g.\ with Control-C). Moreover, if we want to examine \emph{all} pending messages, we would need to call \texttt{receive/1} repeatedly, without necessarily knowing how many messages there are -- which again risks blocking, or that messages are left behind.

A more convenient tool is the utility predicate \texttt{flush/0}. It prints (and removes) all messages currently in the shell's mailbox and, crucially, it never blocks.

Below, we demonstrate both of these utilities by first sending yet another message to our shell and then use \texttt{flush/0} to inspect and empty the contents of its mailbox:

\begin{verbatim}
?- $Self ! goodbye.
true.
?- flush.
Shell got hello
Shell got goodbye
true.
?-
\end{verbatim}

\noindent The \$Var substitution mechanism (or ``dollar notation'') in combination with \texttt{flush/0} comes in very handy during interactive programming in the shell and we shall rely on them extensively when running examples through all of the book.

%Having to copy and paste pids is not very convenient. Instead we can make use of a shell utility feature, borrowed from SWI-Prolog, which allows a variable binding resulting from the successful execution of a shell query to be reused in future shell queries if a dollar symbol is put in front of it, like so:
%
%\begin{verbatim}
%?- $Self ! hello.
%true.
%?-
%\end{verbatim}
%
%\noindent Here, the shell kept track of the most recent binding of \texttt{Self} to a value (\texttt{85234512} in this case) and substituted the term \texttt{\$Self} with that value before the query was run. %\footnote{We emphasize that dollar-prefixed variables are not global variables -- they only work in a terminal attached to a shell.}
%
%While being able to inspect the contents of the shell's mailbox during interactive programming is obviously important, using \texttt{receive/1} is not the most convenient way of doing it. After all, we cannot always be sure that there \textit{are} any messages in the mailbox, and, as noted above, if it is empty then \texttt{receive/1} will block and make normal interaction impossible. The call will just hang indefinitely, and the only way out may be to abort it using Control-C. Furthermore, we may want to see \textit{all} messages in the mailbox which means that we need to run \texttt{receive/1} more than once, but we may not know how \textit{many} messages there are, and then again run into the blocking problem. A more convenient tool here is the utility predicate \texttt{flush/0}. This predicate will allow us to inspect (and remove) all messages from the shell's mailbox and will never block. 
%
%Below, we demonstrate both of these utilities by first sending yet another message to our shell and then use \texttt{flush/0} to inspect and empty the contents of its mailbox:
%
%\begin{verbatim}
%?- $Self ! goodbye.
%true.
%?- flush.
%Shell got hello
%Shell got goodbye
%true.
%?-
%\end{verbatim}
%
%\noindent The \$Var substitution mechanism (or ``dollar notation'') in combination with \texttt{flush/0} comes in very handy during interactive programming in the shell and we shall rely on them extensively when running examples through all of the book.

\subsection{Actors talking to other actors}

It takes two to tango, and an actor talking to itself is not very interesting. Fortunately, as already mentioned, an actor is capable of spawning other actors and then start talking to them. 

The built-in predicate \texttt{spawn/1-3} is used to create new actor processes. In a call such as \texttt{spawn(Goal,\,Pid)} it expects a callable goal to be passed in the first argument, and will bind the variable in the second argument to the pid of the actor created. Here is how this might look like:

\begin{verbatim}
?- spawn(echo_actor, Pid).
Pid = 72347585.
?-
\end{verbatim}

\noindent The goal calls a predicate which determines the behavior of the actor. The actor process runs in parallel with the caller, the shell in this case. Here is an example script that implements a simple echo server:

\begin{verbatim}
echo_actor :-          
    receive({            
        echo(From, Msg) ->
            From ! echo(Msg),   
            echo_actor  
    }).        
\end{verbatim}

\noindent We have seen this piece of code before, and it is time to explain how it works. The predicate \texttt{echo\_actor/0} defines a loop where a call to  \texttt{receive/1} tries to select a message of the form \texttt{echo(From,\,Msg)} from the mailbox and send the echo message back to the actor referenced by the pid to which the variable \texttt{From} is bound, and then continue looping by making a recursive call. %Here is how an actor following this script can be created:

To make our server return the echo message to the shell we must send it a message of the form \texttt{echo(Pid,\,Msg)} with \texttt{From} bound to the pid of the shell and \texttt{Msg} bound to the message. Here is how we do that, this time supported by our shell utilities:

\begin{verbatim}
?- self(Self), $Pid ! echo(Self, hi).
Self = 85234512.
?- flush.
Shell got echo(hi)
true.
?-
\end{verbatim} 

\noindent The sending is asynchronous, i.e. \texttt{!/2} does not block waiting for a response but continues immediately. Also, sending a message to a pid always succeeds and never throws an error, even if the pid points to a non-existing process. 
%Distributed message sending, that is, if \texttt{Actor} has the form \texttt{Pid@Node} or \texttt{Name@Node}, also never throws an exception.
In that case the message is simply discarded. This means that short of having the targeted actor return a confirmation message we will not always know if the message reached it. However, the above case did not leave us in the dark, as the actual echo served as a confirmation message.

Web Prolog guarantees per-sender FIFO delivery: for any fixed sender--receiver pair, messages are enqueued in the receiver's mailbox in the order they were sent. There is no single total order across different senders; concurrent sends to the same receiver may be observed in either order by \texttt{receive/1-2}. 

\subsubsection{Monitoring}

In a call such as \texttt{spawn(Goal,\,Pid,\,Options)} the optional third argument is a list of options used for configuring the actor. One such option is \texttt{monitor}, which if set to \texttt{true} requests that the parent process be notified when the spawned actor terminates.

In the following example, \texttt{monitor(true)} instructs the child actor to send a \texttt{down} message to its parent just before terminating:

\begin{verbatim}
?- spawn(2 > 1, Pid, [
       monitor(true)
   ]).
Pid = 60367387.
?- flush.
Shell got down(60367387,60367387,true)
true.
?-
\end{verbatim}

\noindent A \texttt{down} message has the form \texttt{down(Ref,\,Pid,\,Reason)}
where \texttt{Pid} is the process identifier of the terminated actor, \texttt{Reason} describes \emph{why} it terminated, and \texttt{Ref} identifies the \emph{monitor instance} that produced this notification. When monitoring is installed via \texttt{monitor(true)}, the monitor reference is chosen to be identical to the child's pid, which is why the first two arguments coincide in the example.

%The important point about the \texttt{monitor(true)} option is that it is \emph{atomic} with process creation: the child is either spawned with monitoring already in place, or it is not spawned at all. This avoids a race in which a short-lived child could terminate before monitoring has been established.

In addition to the \texttt{monitor} option, the predicate \texttt{monitor(Pid,\,Ref)} is provided, which installs a monitor and returns a fresh reference in \texttt{Ref}. The reference matters whenever there may be more than one monitor, because each monitor produces its own \texttt{down} message and we must be able to tell them apart:

\begin{verbatim}
?- spawn(sleep(3), Pid),
   monitor(Pid, Ref1),
   monitor(Pid, Ref2).
Pid = 41544343,
Ref1 = 10924110,
Ref2 = 92041933.
?- flush.
Shell got down(10924110,41544343,true)
Shell got down(92041933,41544343,true)
true.
?-
\end{verbatim}

\noindent Here the same parent process installs two monitors on the same child. When the child terminates after three seconds, two \texttt{down} messages are delivered, one per monitor. The distinct references (\texttt{Ref1} and \texttt{Ref2}) identify which monitor instance each message corresponds to.

Using the \texttt{monitor} option is the \emph{safest} way to establish monitoring, because it is \emph{atomic} with respect to process creation: either the child is spawned with monitoring already in place, or it is not spawned at all. In contrast, calling \texttt{monitor/2} after \texttt{spawn/2-3} is a separate step and therefore not atomic. This means that a very short-lived child may terminate in the gap between the \texttt{spawn} returning and the monitor being installed, in which case the parent will never receive a \texttt{down} message. For actors that may complete immediately (or fail immediately), \texttt{monitor(true)} avoids this race condition.

Monitoring can be stopped by passing the monitor ref to \texttt{demonitor/1-2}. When monitoring in installed using the \texttt{monitor} option, the pid of the monitored actor can also be used since it is identical to the ref.

\medskip

\noindent The above example of the use of \texttt{monitor/2} was trivial, but in Chapter~5, both \texttt{monitor/2} and \texttt{demonitor/1-2} will be put to more serious use.



%\noindent In Chapter~5 we return to monitoring in more depth and introduce the predicate \texttt{monitor/2}, which allows monitors to be installed on already running actors and supports multiple independent monitors on the same target process.

%\subsubsection{Monitoring}
%
%An actor can be \emph{monitored}, meaning that another actor will receive a notification when it terminates. Monitoring is \emph{one-way}: the monitored actor is unaffected, while the monitoring actor gains the ability to observe its termination.
%
%\paragraph{Monitoring via \texttt{spawn/3}.}
%The option list in \texttt{spawn(Goal,\,Pid,\,Options)} may include \texttt{monitor(true)}. When this option is used, the parent process will receive a \texttt{down/3} message as the child terminates:
%
%\begin{verbatim}
%?- spawn(2 > 1, Pid, [monitor(true)]).
%Pid = 60367387.
%?- flush.
%Shell got down(60367387,60367387,true)
%true.
%?-
%\end{verbatim}
%
%\noindent A \texttt{down} message has the form \texttt{down(Ref,\,Pid,\,Reason)}
%where \texttt{Pid} is the process identifier of the terminated actor, \texttt{Reason} describes \emph{why} it terminated, and \texttt{Ref} identifies the \emph{monitor instance} that produced this notification.
%
%When monitoring is installed via \texttt{monitor(true)} in \texttt{spawn/3}, the monitor reference is chosen to be identical to the child's pid. This is why the first two arguments coincide in the example above.
%
%\paragraph{Monitoring via \texttt{monitor/2}.}
%In addition to the \texttt{monitor} option, the predicate \texttt{monitor(Pid,\,Ref)} is provided, which installs a monitor and returns a fresh reference in \texttt{Ref}. The reference matters whenever there may be more than one monitor, because each monitor produces its own \texttt{down} message and we must be able to tell them apart:
%
%\begin{verbatim}
%?- spawn(sleep(3), Pid),
%   monitor(Pid, Ref1),
%   monitor(Pid, Ref2).
%Pid = 41544343,
%Ref1 = 10924110,
%Ref2 = 92041933.
%?- flush.
%Shell got down(10924110,41544343,true)
%Shell got down(92041933,41544343,true)
%true.
%?-
%\end{verbatim}
%
%\noindent Here the same parent process installs two monitors on the same child. When the child terminates after three seconds, two \texttt{down} messages are delivered, one per monitor. The distinct references (\texttt{Ref1} and \texttt{Ref2}) identify which monitor instance each message corresponds to.
%
%Using \texttt{monitor(true)} is the \emph{safest} way to establish monitoring, because it is \emph{atomic} with respect to process creation: either the child is spawned with monitoring already in place, or it is not spawned at all. In contrast, calling \texttt{monitor/2} after \texttt{spawn/2-3} is a separate step and therefore not atomic. This means that a very short-lived child may terminate in the gap between the \texttt{spawn} returning and the monitor being installed, in which case the parent will never receive a \texttt{down} message. For actors that may complete immediately (or fail immediately), \texttt{monitor(true)} avoids this race condition.
%
%Monitoring can be stopped using \texttt{demonitor/1-2}, either by providing the monitor reference or (when applicable) the pid of the monitored actor.


%\subsubsection{Monitoring}
%
%In a call such as \texttt{spawn(Goal,\,Pid,\,Options)} the optional third argument is a list of options used for the configuration of the actor. The \texttt{monitor} option can be used to specify what should happen when the actor terminates, i.e what the parent might learn about the reason for its death.
%
%In the following example, the \texttt{monitor} option is set to \texttt{true}, instructing the actor to send a special-purpose \texttt{down} message carrying a small piece of information on how the run went to its parent process just before terminating. The idea is to allow the parent process to observe the child process and to detect if it has terminated for any reason:
%
%\begin{verbatim}
%?- spawn(2 > 1, Pid, [
%       monitor(true)
%   ]).
%Pid = 60367387.
%?- flush.
%Shell got down(60367387,60367387,true)
%true.
%?-
%\end{verbatim}
%
%\noindent In this case, the \texttt{down} message was delived to its parent as soon as the goal \texttt{2 > 1} terminated. The message has the following form, where \texttt{Ref} is a \textit{reference} (a \textit{ref} for short), \texttt{Pid} is a process identifier, and Reason is the reason for terminating:
%
%\begin{verbatim}
%down(Ref, Pid, Reason)
%\end{verbatim}
%
%\noindent Note that in this case, the ref and the pid are identical. This is always the case when the \texttt{monitor} option is used to install a monitor.
%
%Normally, an actor process will terminate and completely disappear when it has run out of things to do. It may be because the script succeeds, or fails, or throws an error. In this case, the atom \texttt{true} in the third argument of the \texttt{down} message tells us that it terminated in a way that can be considered normal.\footnote{Erlang uses the atom \texttt{normal} to indicate this.}
%
%In addition to the \texttt{monitor} option, Web Prolog also supports a built-in predicate \texttt{monitor/2} which takes a pid and an unbound variable that will be bound to a ref. The reference is essential if more than one monitor exists. Otherwise, we cannot tell which monitor triggered the \texttt{down} message. Here is an example demonstrating how it works:
%
%\begin{verbatim}
%?- spawn(sleep(3), Pid),
%   monitor(Pid, Ref1),
%   monitor(Pid, Ref2).
%Pid = 41544343,
%Ref1 = 10924110,
%Ref2 = 92041933.
%?- flush.
%Shell got down(10924110,41544343,true)
%Shell got down(92041933,41544343,true)
%true.
%?-
%\end{verbatim}
%
%\noindent The same parent process sets two monitors on the child, getting two different monitor references (\texttt{Ref1} and \texttt{Ref2}). When the child (running \texttt{sleep(3)}) terminates after three seconds, two \texttt{down} messages reach the shell's mailbox, one per monitor. Each \texttt{down} message includes the unique ref, allowing us to distinguish which monitor instance it corresponds to. Note that when \texttt{monitor/2} is used, the ref is still an integer, but distinct from the pid.
%
%Web Prolog also support a built-in predicate \texttt{demonitor/1-2}, capable of stopping the monitoring of an actor before it delivers a \texttt{down} message. It can be called with either a ref or a pid.

\subsubsection{Linking}

\noindent If the value of the \texttt{link} option is \texttt{true} it means that when an actor terminates, any children that it might have are forced to terminate too. In our proposal, the default value for this option is \texttt{true}. So in the following example, when the outer call to \texttt{spawn/2} terminates, the nonterminating actor spawned by the inner call to \texttt{spawn/1} is automatically killed: 

\begin{verbatim}
?- spawn(spawn(receive({})), Pid).
Pid = 10782012.
?-
\end{verbatim}

\noindent Calling the built-in predicate \texttt{actors/1} confirms that only one actor is alive. We can infer that this must be the shell process, and that the actor created by the call \texttt{spawn(receive(\{\}))} is gone:

\begin{verbatim}
?- actors(Alive).
Alive = [91033267].
?-
\end{verbatim}

\noindent Linking in Web Prolog does \textit{not} mean that a parent actor must terminate if any child that it may have spawned terminates:

\begin{verbatim}
?- spawn((spawn(exit(kill)),receive({})), Pid).
Pid = 40276924.
?-
\end{verbatim}

\noindent Even though the child killed itself immediately, calling \texttt{actors/1} again shows that its parent \texttt{40276924} is still alive:

\begin{verbatim}
?- actors(Alive).
Alive = [91033267, 40276924].
?-
\end{verbatim}

\noindent The link is \textit{unidirectional,} and makes the life of a child dependent on the life of its parent, but never the other way around. (Erlang is different here, as its links are bidirectional. We will discuss this difference and justify our design choice in Section~\ref{sec:discussion-revised}.)

The following example shows that when the \texttt{link} option is set to \texttt{false} the child (created by the inner call to \texttt{spawn/3}) does not necessarily terminate if the parent (created by the outer call to \texttt{spawn/1}) terminates:

\begin{verbatim}
?- spawn(spawn(receive({}),Pid,[link(false)])).
Pid = 10023812.
?- actors(L).
L = [91033267, 24917621].
?-
\end{verbatim}

\noindent In our proposal, the default value for \texttt{link} is \texttt{true}, as this is usually what we want. In fact, and for a reason that will be explained in Chapter~\ref{sec:the-prolog-web}, in the context of distributed web programming it may often be a good idea to \textit{require} that the value of \texttt{link} is set to \texttt{true}, and perhaps only allow the owner of a node to set it to \texttt{false}. 

In later chapters we restrict who may disable lifetime coupling, especially for remote placement, since this interacts with resource control and deployment on the Prolog Web.

\subsubsection{Registering}

\noindent Calling \texttt{register(Name,Pid)} associates the atom \texttt{Name} with \texttt{Pid}. The name can be used instead of the pid when calling \texttt{!/2}. For example, we can register our shell under the name \texttt{shell}:

\begin{verbatim}
?- self(Self), register(shell, Self).
Self = 85234512.
?- shell ! hello.
true.
?- spawn(shell ! goodbye).
true.
?- flush.
Shell got hello
Shell got goodbye
true.
?- 
\end{verbatim}

\noindent Registering is useful when we want to offer a service that should always be available under a name that is easy to remember. Should a crash occur, all our system needs to do is to restart it and associate the same name with the pid of the new process.

\subsubsection{Exiting}

Web Prolog supports two predicates, \texttt{exit/1} and \texttt{exit/2}, that can be used to terminate an actor process. 

If the process calls \texttt{exit(Reason)} it will terminate immediately, and if monitored by its parent process, the parent will be sent a \texttt{down} message, containing the term that \texttt{Reason} was bound to when predicate was called. Here is a simple and somewhat silly example where a process is spawned and monitored by the shell. Since the process is told to exit immediately with the reason \texttt{my\_reason}, the shell receives a \texttt{down} message with the atom \texttt{my\_reason} in its second argument:

\begin{verbatim}
?- spawn(exit(my_reason), Pid, [
       monitor(true)
   ]).
Pid = 91325643.
?- flush.
Shell got down(91325643,91325643, my_reason)
true.
?-
\end{verbatim}

\noindent It sometimes happens that we need to terminate an actor process by force. Our echo server, for example, can only be terminated in this way. The predicate \texttt{exit/2} can be used to terminate any actor process with a known pid. Calling \texttt{exit(Pid,\,Reason)} sends an exit \textit{signal} (not a message) to the actor with that pid, killing it.

If we do not know the pid, but it has a name that we know, we can use that instead. To see how this works, let us spawn a new monitored echo server and register it: 

\begin{verbatim}
?- spawn(echo_actor, Pid, [
       monitor(true)
   ]),
   register(echo_actor, Pid).
Pid = 21562390.
?-
\end{verbatim}

\noindent Now, let us say we change our minds and want to get rid of it again:

\begin{verbatim}   
?- whereis(echo_actor, Pid),
   exit(Pid, 'We changed our minds!').
Pid = 21562390.
?- flush.
Shell got down(21562390,21562390,'We changed our minds!')
true.
?-
\end{verbatim}

\noindent By calling \texttt{exit/2} with the pid in the first argument and a term detailing the reason for exiting in the second, a signal was sent to the process that terminated it. Note that a call to \texttt{exit/2} will only accept a pid in its first argument, so if all we have is a name, the built-in predicate \texttt{whereis/2} must be used to locate it.

Note that the reason passed to \texttt{exit/1-2} can be any term, not just an atom. This means that it can be used to transfer an arbitrarily large chunk of information back to the parent. However, in most cases an atom is all that is needed, and it is worth noticing that using \texttt{true} as a reason can be used to suggest to the parent that the termination was \textit{normal}, even though it was caused by a call to \texttt{exit/1-2}.

What might seem a bit odd is that even thought \texttt{21562390} is now provably dead, trying to send it a message using a pid or calling \texttt{exit/2} still succeeds without an error, although no \texttt{down} message is being sent:

\begin{verbatim}
?- self(Self), $Pid ! echo(Self, bye).
Self = 85234512.
?- exit($Pid, 'We changed our minds!').
true.
?- flush.
true.
?-
\end{verbatim}

\noindent Treating \texttt{!/2} and \texttt{exit/2} as no-ops when the pid points to a non-existent process is consistent with how it works in Erlang. However, trying to send a message using the \textit{name} of a non-existent process generates an error:

\begin{verbatim}
?- echo_actor ! echo($Self, bye).
Error: The name 'echo_actor' is not associated with a pid.
?-
\end{verbatim}

\noindent This too is consistent with how Erlang works.

\subsubsection{An hierarchy of actors}

When an actor process $A_1$ spawns an actor $A_2$, $A_2$ becomes the \textit{child} of $A_1$, and $A_1$ the \textit{parent} of $A_2$. Since $A_2$ may in turn spawn other processes, the actors involved may form a hierarchy. 

\begin{verbatim}
demo :- 
    spawn(parent).

parent :-                         
    spawn(child, Pid, [
        monitor(true)
    ]),          
    receive({                 
        down(_, Pid, Why) ->       
            format("~p ~p~n", [Pid, Why]) ;
        stop -> 
            format("stopped~n")
    }).                       
                             
child :-                         
    spawn(grandchild),              
    receive({
        crash -> 
            exit(crashed) ;
        stop -> 
            format("stopped~n")
    }).                    
                             
grandchild :- 
    receive({stop -> true}). 
\end{verbatim}

\noindent In Erlang, things work a bit differently.

% \vspace{1cm}

% \begin{ErlangBox}{Monitoring is post-hoc here}\label{bx:monitor-posthoc}
% In Erlang, \emph{linking} is a bidirectional relationship between two processes created explicitly (e.g.\ \texttt{spawn\_link/1}) or implicitly (\texttt{link/1}). If one linked process terminates, the other process is immediately notified with an \emph{exit signal}. The default rule is harsh but simple: if a process receives an exit signal and it is not trapping exits, it dies too (with the same reason), which can cause cascades.

% \begin{verbatim}
% demo() -> 
%     spawn(fun parent/0).

% parent() ->                          
%     process_flag(trap_exit,true),
%     Pid = spawn_link(fun child/0), 
%     receive                      
%         {'EXIT', Pid, normal} ->         
%             io:format("~p ~p~n", [Pid, normal]) ;
%         {'EXIT', Pid, Why} ->         
%             io:format("~p ~p~n", [Pid, Why])
%     end.                         
                                
% child() ->                      
%     spawn_link(fun grandchild/0),
%     receive
%         crash -> 
%             exit(crashed) ;
%         stop ->
%             io:format("stopped~n")     
%     end.        
                                
% grandchild() -> 
%     receive stop -> ok end.         
% \end{verbatim}

% \noindent In our example:
% \begin{itemize}
% \item \texttt{p/0} sets \texttt{trap\_exit} and then evaluates \texttt{P\,=\,spawn\_link(fun c/0)}, so \texttt{p} and \texttt{c} are linked.
% \item \texttt{c/0} also \texttt{spawn\_link}s \texttt{gc/0}, so \texttt{c} and \texttt{gc} are linked.
% \item When \texttt{c} receives \texttt{crash}, it calls \texttt{exit(crashed)}. Because \texttt{c} is linked to \texttt{p}, an exit signal with reason \texttt{crashed} is sent to \texttt{p}.
% \item Since \texttt{p} is trapping exits, it does not die; it receives \texttt{\{'EXIT',\,P,\,crashed\}} and prints it.
% \item Separately, because \texttt{c} is linked to \texttt{gc}, \texttt{gc} will also receive an exit signal from \texttt{c} and (since \texttt{gc} is not trapping exits) will typically die as well.
% \end{itemize}

% \end{ErlangBox}

\subsection{A closer look at receive/1-2}

\noindent As shown already, but only for a trivial case, an actor process uses the receive primitive to extract messages from its mailbox. Since the syntax and semantics of \texttt{receive/1-2} is fairly complex, a closer look and more examples are needed. Below we give several simple examples illustrating different ways to use the  \texttt{receive/1-2} predicate in Web Prolog, demonstrating how to handle different \textit{types} of messages, use \textit{timeouts}, apply \textit{guards}, and more. Other examples illustrate how messages are deferred and handled in subsequent \texttt{receive/1-2} calls, demonstrating the flexibility of passing and processing messages in Web Prolog.

\subsubsection{Basic receive}

In Web Prolog, just like in Erlang, the receive operation specifies an ordered sequence of \textit{receive clauses} delimited by semicolons. A receive clause always has a \textit{head} (a term) and a \textit{body} of Prolog goals. Schematically, the basic form of a receive call looks like this:

\begin{verbatim}
receive({  
    Head1 -> Body1 ;
    Head2 -> Body2 ;
    ...
    HeadN -> BodyN
})
\end{verbatim}

\noindent Often, the head is just a single term serving as a \textit{pattern}. Any term will do, ground or non-ground, and even a bare variable is fine.

As in Erlang, \texttt{receive/1} scans the mailbox looking for the first message (i.e. the oldest) that matches a pattern in any of the receive clauses, blocking if no such message is found. If a matching clause is found, the message is removed from the mailbox and the body of the clause is called. In Web Prolog, just like in Erlang, values of any variables bound by the matching of the pattern with a message are available in the body of the clause.

Often, a \texttt{receive/1} call waits for a specific
message and executes the corresponding code in the body of the receive clause if a message that matches the pattern shows up. For example, the following call waits for a message in the form of \texttt{hello(Name)} and prints a greeting when it appears:

\begin{verbatim}
receive({
    hello(Name) ->
        format("Hello, ~s!~n", [Name])
})
\end{verbatim}

\noindent We can specify multiple patterns in a single \texttt{receive/1} call. For example, this call handles messages of either the form \texttt{hello(Name)} or the form \texttt{goodbye(Name)}, but only one of them:

\begin{verbatim}
receive({
    hello(Name) ->
        format("Hello, ~s!~n", [Name]) ;
    goodbye(Name) ->
        format("Goodbye, ~s!~n", [Name])
})
\end{verbatim}

\noindent We can use the \texttt{\_} pattern to catch any message. In the following case, any message that does not match \texttt{hello(Name)} will be caught by the \texttt{\_} pattern:

\begin{verbatim}
receive({
    hello(Name) ->
        format("Hello, ~s!~n", [Name]) ;
    _ ->
        format("Unknown message received.~n")
})
\end{verbatim}

\noindent We can nest \texttt{receive/1-2} calls. Here, after having received the \texttt{start(Name)} message, the call waits for a \texttt{continue(Msg)} message:

\begin{verbatim}
receive({
    start(Name) ->
        format("Starting with ~s.~n", [Name]),
        receive({
            session(Msg) ->
                format("Continuing with ~s~n", [Msg])
        })
})
\end{verbatim}

\noindent 

\subsubsection{Messages deferred}

If no pattern matches a message in the mailbox, the message is \textit{deferred}, which means that the message does not match any of the  patterns in the current \texttt{receive/1-2} call and remains in the process's mailbox, possibly to be handled later in the control flow of the process. The receive is still running, waiting for more messages to arrive, and for one that will match. Some simple examples illustrating this behavior are shown below.

In the following example, the \texttt{goodbye("Bob")} message -- which is the oldest and therefore first in line -- does not match the clause in the first receive call and is deferred. The \texttt{hello("Alice")} message matches that clause, and then the clause in the second receive call handles the \texttt{goodbye("Bob")} message:

\begin{verbatim}
?- self(Self),
   Self ! goodbye("Bob"),
   Self ! hello("Alice"),
   receive({
       hello(Name1) ->
           format("Hello, ~s!~n", [Name1])
   }),
   receive({
       goodbye(Name2) ->
           format("Goodbye, ~s!~n", [Name2])
   }).
Hello, Alice!
Goodbye, Bob!
true.
?-
\end{verbatim}

\noindent This behavior is particularly useful if we expect two messages but are not sure which one will arrive first. For example, if we insist on processing a message \texttt{foo} before \texttt{bar}, we can easily do that with receive, like so:

\begin{verbatim}
wait_foo :-                   wait_bar :-
    receive({                     receive({
        foo ->                        bar ->
            process_foo,                  process_bar
            wait_bar              }).                             
    }).
\end{verbatim}

\noindent Even if \texttt{bar} arrives in the mailbox before \texttt{foo}, calling \texttt{wait\_foo/0} would result in \texttt{foo} being selected and processed before \texttt{bar}. This is why the receive operator is often referred to as \textit{selective} receive. 

\subsubsection{Guards}\label{sec:receive-guards}

The head of a receive clause can, in addition to the pattern, optionally specify a \textit{guard} in the form of a Prolog goal. The role of a head of  the form \texttt{Pattern\,if\,Goal} is to make pattern matching more expressive. Here is an example that distinguishes between positive and non-positive numbers using guards:

\begin{verbatim}
receive({
    number(N) if N > 0 ->
        format("Positive number: ~p~n", [N]) ;
    number(N) if N =< 0 ->
        format("Non-positive number: ~p~n", [N])
})
\end{verbatim}

\noindent That was simple, and will work in Erlang too, but here is another example, using a different \textit{kind} of goal after the \texttt{if} operator:

\begin{verbatim}
?- self(Self),
   Self ! hello(xantippa),
   receive({
       hello(W) if husband(W, H) ->
           format("Hello, ~w, say hello to ~w!~n", [W,H])
   }).
Hello, xantippa, say hello to socrates!
true.
?-
\end{verbatim}

\noindent Note that the variable \texttt{H} does not occur in the pattern. A guard like that cannot be used in Erlang. In Web Prolog, its value can be passed to the body of the clause and do a job there. So while readers familiar with Erlang may wonder why we choose \texttt{if} instead of \texttt{when}, which is the operator Erlang is using, we just happen to think that this difference is significant enough to warrant a different name for the operator.


\subsubsection{Timeouts}

\noindent The optional second argument of \texttt{receive/1-2} expects a list of options. The \texttt{timeout} option takes an integer or float that specifies the number of seconds before the call will succeed anyway, even if no match has been found. The \texttt{on\_timeout} option takes a goal that is called if timeout occurs. Here is an example of its use:

\begin{verbatim}
receive({
    hello(Name) ->
        format("Hello, ~s!~n", [Name])
},[ timeout(5),
    on_timeout(format("No match received in 5 seconds.~n"))
])
\end{verbatim}

\noindent If no message of the form \texttt{hello(Name)} is received within 5 seconds, the code in the \texttt{on\_timeout} option is executed.


\subsubsection{An implementation of sleep/1}

Here is how a predicate \texttt{sleep/1} that suspends execution \texttt{Time} seconds can be defined:

\begin{verbatim}
sleep(Time) :-
    receive({},[ 
        timeout(Time)
    ]).
\end{verbatim}

\noindent The call to \texttt{receive/2} does not wait for messages, but when \texttt{Time} seconds have passed, it still succeeds.

\subsubsection{An implementation of flush/0}

As we noted earlier, being able to inspect the contents of the shell's mailbox during interactive programming is important, and \texttt{flush/0} is a nice tool for doing just that. Its definition also serves as yet another example of the use of the \texttt{timeout} option:

\begin{verbatim}
flush :-
    receive({
        Message -> 
            format("Shell got ~q~n",[Message]),
            flush
    },[ 
        timeout(0)
    ]).
\end{verbatim}

\noindent The value \texttt{0} of the \texttt{timeout} option ensures that the loop terminates immediately if no messages remain in the mailbox. This is how the hanging of the \texttt{receive/1} call is avoided.

\subsubsection{A simple priority queue}\label{sec:priority-queue}

To demonstrate the use of the \texttt{if} operator and the use of two \texttt{receive/2} options that causes a goal to run on timeout, we show a priority queue example borrowed from Fred H\'ebert's textbook on Erlang \cite{Hebert:2013:LYE:2543986}. The purpose is to build a list of messages with those with a priority above 10 coming first:

\begin{verbatim}
important(Messages) :-
    receive({
        Priority-Message if Priority > 10 ->
            Messages = [Message|MoreMessages],
            important(MoreMessages)
    },[ timeout(0),
        on_timeout(normal(Messages))
    ]).

normal(Messages) :-
    receive({
        _-Message ->
            Messages = [Message|MoreMessages],
            normal(MoreMessages)
    },[ timeout(0),
        on_timeout(Messages=[])
    ]).
\end{verbatim}

\noindent The timeout set to \texttt{0} means that it will occur immediately, but the system tries all messages currently in the mailbox first.

Below, we test this program by first sending four messages to the shell process, and then calling \texttt{important/1}:

\begin{verbatim}
?- self(S),
   S ! 15-high, S ! 7-low, S ! 1-low, S ! 17-high.
S = 34871244.
?- important(Messages).
Messages = [high,high,low,low].
?- 
\end{verbatim}

\noindent Note that the implementation of the priority-queue example relies on the deferring behavior of \texttt{receive/1-2} and would not work without it.

\subsubsection{The receive predicate is semi-deterministic}

The predicate \texttt{receive/1-2} is \textit{semi-deterministic}, i.e. it either fails, or succeeds exactly once. The only way it will fail is if the goal in the \textit{body} of one of its receive clauses fails, or if timeout occurs and the goal passed in the \texttt{on\_timeout} option fails. 

To see how it pans out in a simple corner case, consider the following two calls:

\begin{verbatim}
receive({foo(X) -> true})       receive({foo(X) -> fail})
\end{verbatim} 

\noindent The first call will succeed if a message matching the pattern \texttt{foo(X)} appears in the mailbox of the actor process executing the call, a term such as \texttt{foo(314)} for example. The second call will fail (and possibly cause backtracking) once \texttt{foo(314)} appears. Only by the first call will the variable \texttt{X} be bound (to \texttt{314}). Both calls will remove the matched message from the mailbox. In both cases, if a message appears that does not match the pattern, it is deferred.

To implement a looping behavior, Prolog programmers occasionally use a \textit{failure-driven} loop that relies on backtracking rather than recursion. Using this technique, our echo server can be rewritten like so:

\begin{verbatim}
echo_actor  :-
    repeat,        
    receive({             
        echo(From, Msg) -> 
            From ! echo(Msg),    
            fail   
    }).
\end{verbatim}

\noindent For an Erlang programmer, this use of \texttt{receive/1} may come as a surprise and is not a technique that can be used in Erlang. In this particular case, a Web Prolog programmer would be advised to stick to the recursive version. However, there are cases when a failure-driven loop is the only way forward and we shall look at an important such case towards the end of this chapter. %Chapter~\ref{sec:prolog-agents}.



\section{Erlang-style programming examples in Web Prolog}

\begin{quotation}
\noindent Reading the code was fun -- I had to do a double take -- was I reading Erlang or Prolog -- they often look pretty much the same.
\vspace{-0.4cm}
\flushright\textit{Joe Armstrong} (p.c. June 18, 2018)
\vspace{0.1cm}
\end{quotation}

\noindent Not only do Web Prolog programs \textit{look} a lot like Erlang, they \textit{behave} a lot like Erlang too. A good way to demonstrate this is to translate a fair number of different Erlang programming examples borrowed from tutorials and text books into Web Prolog and show that they run just like in Erlang. In this section we shall look at a count server, a fridge simulation, a universal stateful server with hot code swapping, actors playing ping-pong, and an approach to building simple supervision hierarchies. They are all concurrent programs, but only locally so. (Distributed programming will be dealt with in Chapter~\ref{sec:the-prolog-web}.)

During the exploration of the examples, we point to the importance of the use of send and receive for implementing the \textit{communication protocols} allowing actors acting as clients to, still within the bounds of one node, talk to actors acting as servers. We also explain why it is often a good idea to hide the details of such protocols from programmers behind a dedicated predicate API.

For more exhaustive documentation of all built-in predicates available in Web Prolog, consult Appendix~\ref{sec:manual}.

\subsection{A count server}\label{sec:servers}

We have already looked at an example of a \textit{stateless} server, namely the \texttt{echo\_actor} presented earlier in this chapter. Let us now turn to \textit{stateful} servers and show how state may be programmed. In the following example, we demonstrate how to script an actor that can keep a \textit{count}:

\begin{verbatim}
count_actor(Count0) :-                            
    receive({                         
        count(From) ->
            Count is Count0 + 1,              
            From ! count(Count),              
            count_actor(Count) ;
        stop ->
            true
    }).                            
\end{verbatim}

\noindent The predicate \texttt{count\_actor/1} defines a loop where a call to the built-in predicate \texttt{receive/1} tries to select a message of the form \texttt{count(From)} from the mailbox, increment the counter, send the current count back to the actor referenced by the pid to which the variable \texttt{From} is bound, and continue looping by making a recursive call. Note how the state of the counter -- the current count, that is -- is kept in the argument of the predicate. Note also that we added a second receive clause that will allow us to terminate the server without using \texttt{exit/2}. We just have to send it a message of the form \texttt{stop}. 

Here is how an actor following this script can be made to perform when spawning it:

\begin{verbatim}
?- spawn(count_actor(0), Pid, [
       monitor(true)
   ]).
Pid = 60367387.
?-
\end{verbatim}

\noindent The \texttt{monitor} option was set to \texttt{true}, instructing the actor to send a special-purpose \texttt{down} message carrying a small piece of information on how the run went to its parent process just before terminating. %Default is to not monitor.

Calling \texttt{self/1} determines the identity of the toplevel process -- the process that just became the parent of the spawned actor:\footnote{If you are an actor, this is how you find out who you are!}

\begin{verbatim}
?- self(Self).
Self = 41167597.
?-
\end{verbatim}

\noindent In the next step the send operator \texttt{!/2} is used for sending a message to the spawned actor, instructing it to increment the count and to return the result in the form of a message:

\begin{verbatim}
?- $Pid ! count($Self).
true.
?-
\end{verbatim}

%This serves to demonstrate an important principle: If you want an answer, you must provide a return address!
\noindent A call to \texttt{receive/1} can be made in order to collect the message arriving from the actor:

\begin{verbatim}
?- receive({Count -> true}).
Count = count(1).
?-
\end{verbatim}

\noindent Here, \texttt{!/2} is used to send two messages to the actor that will end up in its mailbox, in the order they were sent:

\begin{verbatim}
?- $Pid ! count($Self), $Pid ! stop.
true.
?-
\end{verbatim}

\noindent Finally, the utility predicate \texttt{flush/0} is used to inspect the contents of the mailbox belonging to the toplevel process:

\begin{verbatim}
?- flush.
Shell got count(2)
Shell got down(60367387,60367387, true)
true.
?-
\end{verbatim}

\noindent Because the actor was monitored, a \texttt{down} message was found in addition to the current count. The value \texttt{true} in the second argument means that \texttt{count\_actor/1} succeeded.

\subsection{A bigger, tastier example}

As a tastier example of how a process can be made to hold an updatable state during a conversation we have adapted a fridge simulation example from Fred H\'ebert's introduction to Erlang \cite{Hebert:2013:LYE:2543986}:\footnote{\url{http://learnyousomeerlang.com/more-on-multiprocessing\#state-your-state}}

\begin{verbatim}
fridge(FoodList0) :-
    receive({
        store(From, Food) ->
            self(Self),
            From ! Self-ok,
            fridge([Food|FoodList0]);
        take(From, Food) ->
            self(Self),
            (   select(Food, FoodList0, FoodList)
            ->  From ! Self-ok(Food),
                fridge(FoodList)
            ;   From ! Self-not_found,
                fridge(FoodList0)
            );
        terminate ->
            true
    }).
\end{verbatim}

\noindent The program creates a process allowing three operations: storing food in the fridge, taking food from the fridge, and terminating the fridge. It is only possible to take food that has been stored beforehand. Again, with the help of recursion the state of a process can be held entirely in the argument of the predicate. In this case we choose to store all the food as a list, and then look in that list when someone needs something.

Assuming the above program is already available, the following session creates the server process. We can then call \texttt{self/1}, \texttt{!/2} and \texttt{flush/0} from the shell in order to simulate the actions of a client:

\begin{verbatim}
?- spawn(fridge([]), Pid, [
       monitor(true)
   ]).
Pid = 77346122.
?- self(Me), $Pid ! store(Me, meat), $Pid ! store(Me, cheese).
Me = 97216744.
?- flush.
Shell got 77346122-ok
Shell got 77346122-ok
true.
?- self(Me), $Pid ! take(Me, cheese).
Me = 97216744.
?- flush.
Shell got 77346122-ok(cheese)
true.
?- $Pid ! terminate.
true.
?- flush.
Shell got down(77346122,77346122,true)
true.
?-
\end{verbatim}

\subsection{Hiding the details of protocols}\label{sec:hiding-protocol}

\noindent In the previous example, we expected programmers to know the details of the protocol that must be followed when interacting with our fridge simulation, which forced them to make raw calls using the send operator in combination with the \texttt{flush/0} utility predicate. As suggested by Fred H\'ebert in his book, that is often a useless burden, and good way around it is to abstract messages away with the help of predicates (or in H\'ebert's case, functions) dealing with receiving and sending them:

\begin{verbatim}
store(Pid, Food, Response) :-
    self(Self),
    Pid ! store(Self, Food),
    receive({
        Pid-Response -> true
    }).

take(Pid, Food, Response) :-
    self(Self),
    Pid ! take(Self, Food),
    receive({
        Pid-Response -> true
    }).
\end{verbatim}

\noindent Calling \texttt{receive/1} immediately after having sent the message guarantees that the communication between client and server stays synchronous.  Using \texttt{store/3} and \texttt{take/3}, the interaction with the fridge becomes somewhat easier:

\begin{verbatim}
?- spawn(fridge([]), Pid, [
       monitor(true)
   ]).
Pid = 55289322.
?- store($Pid, cheese, Response).
Response = ok.
?- take($Pid, cheese, Response).
Response = ok(cheese).
?-
\end{verbatim}

\noindent When dealing with actors with more complex protocols, such abstractions turns out to be of considerable value. 

\subsection{Using delayed sending}\label{sec:delayed-sending-impl}

Appendix~\ref{sec:manual} documents a built-in predicate \texttt{send/2-3} that allows a process to delay the sending of a message, but also to cancel the sending if needed, by calling \texttt{cancel/1}. This can be used to implement an alarm that can be canceled, like so:

\begin{verbatim}
alarm :-
    receive({
        ring ->
            writeln('Alarm ringing!'),
            alarm;
        stop ->
            true
    }).
\end{verbatim}

\noindent Here we schedule the alarm to ring in 5 seconds, with an \texttt{id} set to \texttt{alarm1}. Then we cancel the alarm after a couple of seconds.

\begin{verbatim}
?- spawn(alarm, Pid).
Pid = 60034123.
?- send($Pid, ring, [
       delay(5),
       id(alarm1)
   ]).
?- cancel(alarm1).
true.
?-
\end{verbatim}

\noindent If the cancelation succeeds, the message is withdrawn and the alarm never rings. If \texttt{cancel/1} is called too late, the alarm message has already been delivered and ``Alarm ringing!'' will be printed.

\subsection{Prolog actors playing ping-pong}\label{sec:actors-and-concurrency}

\noindent Since the servers in the previous sections are running in parallel to the shell and are talking to it using asynchronous messaging, we have already demonstrated the use of concurrency. Below, in a probably more convincing example inspired by a user's guide to Erlang,\footnote{See \url{http://erlang.org/doc/getting\_started/conc\_prog.html}} two processes are first created and then start sending messages to each other a specified number of times:

\begin{verbatim}
ping(0, Pong_Pid) :-
    Pong_Pid ! finished,
    format('Ping finished.~n',[]).
ping(N, Pong_Pid) :-
    self(Self),
    Pong_Pid ! ping(Self),
    receive({
        pong -> 
            format('Ping received pong.~n',[])
    }),
    N1 is N - 1,
    ping(N1, Pong_Pid).

pong :-
    receive({
        finished ->
            format('Pong finished.~n',[]);
        ping(Ping_Pid) ->
            format('Pong received ping.~n',[]),
            Ping_Pid ! pong,
            pong
    }).

ping_pong :-
    spawn(pong, Pong_Pid),
    spawn(ping(3, Pong_Pid)).
\end{verbatim}

\noindent When \texttt{ping\_pong/0} is called the behavior of this program exactly mirrors the behavior of the original version in Erlang:

\begin{verbatim}
?- ping_pong.
true.
Pong received ping.
Ping received pong.
Pong received ping.
Ping received pong.
Pong received ping.
Ping received pong.
Ping finished.
Pong finished.
?-
\end{verbatim}

\noindent This is a concurrent program, but its execution is not done in parallel, but rather as is illustrated in Figure~\ref{fig:ping-pong-tasks}. 

\begin{figure}[h]
    \centering
	\includegraphics[width=10cm]{ping-pong-tasks}
    \caption{A timeline where the upper line represent the ``pinger'' and the lower line the ``ponger,'' and where each gray bar represent a task consisting of the sending of a message (in light gray) and the reception of the response (in darker gray)}
    \label{fig:ping-pong-tasks}
\end{figure}

\noindent In Appendix~\ref{sec:benchmarking} a slightly modified version of this program is used for benchmarking send and receive in Web Prolog. It turns out that if we run this on a 2023 Apple iMac with the M3 chip with \texttt{N} set to \texttt{100,000} instead of \texttt{3} it runs in less than one second. This means that each task is performed in less than 5 microseconds. Appendix~\ref{sec:benchmarking} also shows that Erlang is even faster.

%\subsection{Event-driven state machines}\label{sec:implementing-state-machines}
%
%Here we highlight a capability of Web Prolog that Erlang lacks and show what it buys us in practice. Guards in receive clauses may be arbitrary Prolog goals whose bindings carry into the clause body; we illustrate this with an event-driven state machine whose guarded transitions can consult the process's knowledge base. 
%
%---
%
%Web Prolog's receive primitive can be used to implement simple event-driven state machines in straightforward code. Patterns in receive clauses can be used to match against event messages, guards to encode conditions, and the bodies of receive clause to perform actions. 
%
%Figure~\ref{fig:state-machine} depicts a machine comprising two states and four transitions.
%
%\begin{figure}[h]
%    \centering
%	\includegraphics[width=9cm]{state-machine}
%    \caption{A simple event-driven state machine.}
%    \label{fig:state-machine}
%\end{figure}
%
%\noindent Using predicate names to encode states, here is how it can be implemented:
%
%\begin{verbatim}
%s0 :-                            s1 :-
%    receive({                        receive({
%       e1(X) if c1(X,Y) ->              e3(X) if c3(X,Y) ->
%           a1(Y),                           a3(Y),
%           s1 ;                             s1 ;
%       e2(X) if c2(X,Y) ->              _AnyEvent  ->
%           a2(Y),                           s0
%           s0                        }).                          
%    }).
%\end{verbatim}
%
%\noindent In state \texttt{s0} the machine is waiting for an event message of the form \texttt{e1(X)} or \texttt{e2(X)} to appear in the mailbox of the current process. If \texttt{e1(X)} is matched and the condition \texttt{c1(X,Y)} holds, the action \texttt{a1(Y)} will be called and the machine will transition to \texttt{s1}, but if \texttt{e2(X)} is matched and \texttt{c2(X,Y)} holds, then \texttt{a2(Y)} is called and the machine stays in state \texttt{s0}. And so on.
%
%As far as we aware, event-driven state machines have not been used much in the Prolog world. Perhaps the reason is that primitives for sending and receiving messages did not appear in Prolog until (a few) platforms implemented the ISO Prolog Threads draft standard.\footnote{\url{http://logtalk.org/plstd/threads.pdf}} As can be seen by an Erlang/OTP behavior such as \texttt{gen\_statem},\footnote{\url{http://erlang.org/doc/man/gen\_statem.html}} Erlang takes event-driven state machines very seriously, and perhaps it is time for Prolog to follow suit? We shall assume that it is, and in Chapter~\ref{sec:web-prolog-statechart-behavior} present a proposal for how to introduce Web Prolog as a scripting language for StateChart XML, a W3C standard which provides an XML-based notation for event-driven state machines of a particularly sophisticated kind.

\subsection{Event-driven state machines}\label{sec:implementing-state-machines}

Building on the guard mechanism introduced in Section~\ref{sec:receive-guards}, we now show a simple but representative application: an event-driven state machine whose transitions can consult the process's knowledge base. The key point is that clause selection can depend not only on the incoming message pattern, but also on arbitrary relational conditions, and any bindings produced by those conditions can be used by the transition action.

Web Prolog's \texttt{receive} primitive can implement such machines directly. Message patterns represent events, guards represent transition conditions, and clause bodies implement actions and state updates. Figure~\ref{fig:state-machine} depicts a machine with two states and four transitions.

\begin{figure}[h]
    \centering
    \includegraphics[width=9cm]{state-machine}
    \caption{A simple event-driven state machine.}
    \label{fig:state-machine}
\end{figure}

\noindent Using predicate names to encode states, one possible implementation is:

\begin{verbatim}
s0 :-                            s1 :-
    receive({                        receive({
       e1(X) if c1(X,Y) ->              e3(X) if c3(X,Y) ->
           a1(Y),                           a3(Y),
           s1 ;                             s1 ;
       e2(X) if c2(X,Y) ->              _AnyEvent  ->
           a2(Y),                           s0
           s0                        }).                          
    }).
\end{verbatim}

\noindent In state \texttt{s0} the machine waits for an event message of the form \texttt{e1(X)} or \texttt{e2(X)} in the current process's mailbox. If \texttt{e1(X)} matches and \texttt{c1(X,Y)} succeeds, the action \texttt{a1(Y)} is executed and the machine transitions to \texttt{s1}. If \texttt{e2(X)} matches and \texttt{c2(X,Y)} succeeds, the action \texttt{a2(Y)} is executed and the machine remains in \texttt{s0}. The other state is analogous: \texttt{e3(X)} triggers a guarded self-loop in \texttt{s1}, while any other event falls back to \texttt{s0}.

Seen from an Erlang perspective, the example also clarifies the design choice: Erlang restricts guards to a small, carefully controlled language, whereas Web Prolog permits relational guards that consult application predicates and introduce new bindings, as in the \texttt{husband/2} example earlier in this chapter. This expressiveness is exactly what makes knowledge-base-dependent transitions feel natural, but it also places responsibility on the programmer to keep guards efficient and free of side effects.

It is worth noting that event-driven state machines have historically been less prominent in the Prolog world. One likely reason is that message passing arrived relatively late and unevenly across Prolog systems, gaining wider traction only with thread and message-passing facilities such as those explored in the ISO Prolog Threads draft.\footnote{\url{http://logtalk.org/plstd/threads.pdf}} By contrast, Erlang/OTP treats this programming model as foundational, as exemplified by \texttt{gen\_statem}.\footnote{\url{http://erlang.org/doc/man/gen\_statem.html}} In Chapter~\ref{sec:web-prolog-statechart-behavior} we build on this observation and propose Web Prolog as a scripting language for StateChart XML (SCXML), a W3C standard for event-driven state machines of a more sophisticated kind.

===

Erlang enforces purity and efficiency by only allowing a restricted set of primitives in guards, and completely disallows calling user-defined functions. For the use cases the people behind Erlang had, it probably made sense to impose such restrictions. 

In Web Prolog, however, although only the first solution will be searched for, \textit{any} goal may be used as a guard and values of variables bound by it are available in the body. This makes the receive construct more powerful than in Erlang, but it also means that the programmer is made responsible for keeping guards as simple and efficient as possible and to avoid side effects. In our opinion, and as exemplified by the state-machine code in Section~\ref{sec:implementing-state-machines}, enabling the programmer to condition the matching of a receive clause on the content of the entire Prolog database makes it worth it. 

Most Erlang developers accept that guards serve a specific role, and they understand the trade-offs. But some find it limiting and express desire for more expressive, reusable guard logic. So there is broad interest and research into relaxing the restriction for \textit{pure} functions.

---

The change of operator names from \texttt{when} to \texttt{if} suggests the change of semantics. (This is not the only reason why we prefer \texttt{if}.) 

%\subsection{Event-driven state machines}\label{sec:implementing-state-machines}
%
%This section highlights a capability of Web Prolog that Erlang deliberately restricts. In Web Prolog, the guard of a \texttt{receive} clause may be an arbitrary Prolog goal; the first solution is used, and any variable bindings produced by the guard are available in the clause body. This makes it natural to express event-driven state machines whose transitions can consult and exploit the process's knowledge base.
%
%Web Prolog's \texttt{receive} primitive can implement simple event-driven state machines directly: message patterns represent events, guards represent transition conditions, and clause bodies implement actions and state updates.
%
%Figure~\ref{fig:state-machine} depicts a machine with two states and four transitions.
%
%\begin{figure}[h]
%    \centering
%    \includegraphics[width=9cm]{state-machine}
%    \caption{A simple event-driven state machine.}
%    \label{fig:state-machine}
%\end{figure}
%
%\noindent Using predicate names to encode states, one possible implementation is:
%
%\begin{verbatim}
%s0 :-                            s1 :-
%    receive({                        receive({
%       e1(X) if c1(X,Y) ->              e3(X) if c3(X,Y) ->
%           a1(Y),                           a3(Y),
%           s1 ;                             s1 ;
%       e2(X) if c2(X,Y) ->              _AnyEvent  ->
%           a2(Y),                           s0
%           s0                        }).                          
%    }).
%\end{verbatim}
%
%\noindent In state \texttt{s0} the machine waits for an event message of the form \texttt{e1(X)} or \texttt{e2(X)} in the current process's mailbox. If \texttt{e1(X)} matches and the guard \texttt{c1(X,Y)} succeeds, the action \texttt{a1(Y)} is executed and the machine transitions to \texttt{s1}. If instead \texttt{e2(X)} matches and \texttt{c2(X,Y)} succeeds, the action \texttt{a2(Y)} is executed and the machine remains in \texttt{s0}. The other state is analogous: \texttt{e3(X)} triggers a guarded self-loop in \texttt{s1}, while any other event falls back to \texttt{s0}.
%
%This style depends on guard expressiveness. Erlang enforces purity and predictability by allowing only a restricted set of guard expressions and by disallowing calls to user-defined functions inside guards. Given Erlang's design goals, these restrictions are understandable. Web Prolog makes a different trade-off: any goal may appear in a guard, and the bindings it produces may be used in the body. The benefit is that clause selection can be conditioned on the content of the process's knowledge base, not merely on the shape of the incoming message. The cost is that the programmer must keep guards simple and efficient, and avoid side effects, since guards run as part of message selection.
%
%It is worth noting that event-driven state machines have historically been less prominent in the Prolog world. One likely reason is that message passing arrived relatively late and unevenly across Prolog systems, gaining wider traction only with thread and message-passing facilities such as those explored in the ISO Prolog Threads draft.\footnote{\url{http://logtalk.org/plstd/threads.pdf}} By contrast, Erlang/OTP treats this programming model as foundational, as exemplified by \texttt{gen\_statem}.\footnote{\url{http://erlang.org/doc/man/gen\_statem.html}} In Chapter~\ref{sec:web-prolog-statechart-behavior} we build on this observation and propose Web Prolog as a scripting language for StateChart XML (SCXML), a W3C standard for event-driven state machines of a more sophisticated kind.


%\noindent In state \texttt{s0} the machine is waiting for an event message of the form \texttt{e1(X)} or \texttt{e2(X)} to appear in the mailbox of the current process. If \texttt{e1(X)} is matched and the condition \texttt{c1(X,Y)} holds, the action \texttt{a1(Y)} will be called and the machine will transition to \texttt{s1}, but if \texttt{e2(X)} is matched and \texttt{c2(X,Y)} holds, then \texttt{a2(Y)} is called and the machine stays in state \texttt{s0}. And so on.
%
%As far as we aware, event-driven state machines have not been used much in the Prolog world. Perhaps the reason is that primitives for sending and receiving messages did not appear in Prolog until (a few) platforms implemented the ISO Prolog Threads draft standard.\footnote{\url{http://logtalk.org/plstd/threads.pdf}} As can be seen by an Erlang/OTP behavior such as \texttt{gen\_statem},\footnote{\url{http://erlang.org/doc/man/gen\_statem.html}} Erlang takes event-driven state machines very seriously, and perhaps it is time for Prolog to follow suit? We shall assume that it is, and in Chapter~\ref{sec:scxml} present a proposal for how to introduce Web Prolog as a scripting language for StateChart XML, a W3C standard which provides an XML-based notation for event-driven state machines of a particularly sophisticated kind.
%
%---
%
%Erlang enforces purity and efficiency by only allowing a restricted set of primitives in guards, and completely disallows calling user-defined functions. For the use cases the people behind Erlang had, it probably made sense to impose such restrictions. 
%
%In Web Prolog, however, although only the first solution will be searched for, \textit{any} goal may be used as a guard and values of variables bound by it are available in the body. This makes the receive construct more powerful than in Erlang, but it also means that the programmer is made responsible for keeping guards as simple and efficient as possible and to avoid side effects. In our opinion, and as exemplified by the state-machine code in Section~\ref{sec:implementing-state-machines}, enabling the programmer to condition the matching of a receive clause on the content of the entire Prolog database makes it worth it. 
%
%Most Erlang developers accept that guards serve a specific role, and they understand the trade-offs. But some find it limiting and express desire for more expressive, reusable guard logic. So there is broad interest and research into relaxing the restriction for \textit{pure} functions.
%
%---
%
%The change of operator names from \texttt{when} to \texttt{if} suggests the change of semantics. (This is not the only reason why we prefer \texttt{if}.) 
%

\subsection{Getting answers through backtracking}\label{sec:answers-through-backtracking}

We now turn to another feature that Web Prolog has but that Erlang lacks, and give an example that rely on semi-deterministic message reception -- a receive either fails, or succeeds exactly once -- to structure clean, failure-driven loops that interact predictably with asynchronous messaging, including stepping through solution streams and a small toplevel-style interaction.

As stated previously, at least in theory, unexpected interactions between language features and possible impedance mismatches between Prolog's relational, nondeterministic programming model and Erlang's functional and message passing model should not cause any problems. As we now turn to more examples that go beyond what Erlang can easily do, we need to see how well the Erlang-style constructs do mix with backtracking for example? In this section we show some examples suggesting that the mix is both sound and easy to understand. 

Suppose the query given in the argument to \texttt{spawn/2} has several answers, a query such as \texttt{?-human(Who)} for example. Below, a goal containing this query is called, the first solution is sent back to the calling process, and \texttt{receive/1} is then used in order to listen for a message of the form \texttt{next} or \texttt{stop} before terminating:

\begin{verbatim}
?- self(Self),
   spawn(( human(Who),
           Self ! Who,
           receive({
               next -> fail ;
               stop -> true
           })
         ), Pid).
Pid = 76123351,
Self = 90054377.
?- flush.
Shell got socrates
true.
?- $Pid ! next.
true.
?- flush.
Shell got plato
true.
?- $Pid ! stop.
true.
?-
\end{verbatim}

\noindent As this session illustrates, the spawned goal generated the solution \texttt{socrates}, sent it to the mailbox of the parent shell process, and then suspended and waited for more messages. When the message \texttt{next} arrived, the forced failure triggered backtracking which generated and sent \texttt{plato} to the mailbox of the toplevel shell process. The next message was \texttt{stop}, so the spawned process terminated. 
Note that the example demonstrated an actor adhering to what might be seen as a tiny communication protocol accepting only the messages \texttt{next} and \texttt{stop}. 

Looking at the code in the first argument of the calls to \texttt{spawn/2} above, this is how we in Prolog often loop over the solutions to a query, using a failure-driven loop rather than a recursive one. Again, the most obvious way to make the code work as expected, is to allow receive to fail. 

One needs to observe, however, that the goal to be solved in the above example is hard-coded into the program, that the protocol for the communication between client and actor is overly simplistic, and that neither failure of the spawned goal, nor error thrown by it, are handled. There is clearly a need for something more complete and more generic.

\subsection{The toplevel behavior: a preliminary sketch}\label{sec:simple-toplevel}

In essence, a Prolog toplevel is a failure-driven loop running inside another loop -- recursive or failure driven. The failure-driven inner loop allows a client to ask for one solution at a time to a goal, while the outer loop allows it to call several goals in a row and run each of them to completion.

Below, we show how we can build a simple Prolog toplevel by using a meta-predicate (such as \texttt{call\_cleanup/2}) and by specifying a small set of custom messages carrying answers and/or the state of the process that needs to be returned to the calling process. The predicate \texttt{call\_cleanup/2} is here used not only to call a goal, but also to check if any choice points remain after the goal has been called or backtracked into.\footnote{\url{http://www.swi-prolog.org/pldoc/man?predicate=call\_cleanup/2}}

\begin{verbatim}
toplevel(Pid) :-
    toplevel(Pid, []).

toplevel(Pid, Options) :-
    self(Self),
    spawn(session(Pid, Self), Pid, Options).

session(Pid, Parent) :-
    receive({
        '$call'(Template, Goal) ->
            (   call_cleanup(Goal, Det=true), 
                (   var(Det)
                ->  Parent ! success(Pid, Template, true),
                    receive({
                        '$next' -> fail ;
                        '$stop' -> true
                    }),
                    !
                ;   Parent ! success(Pid, Template, false)
                )
            ;   Parent ! failure(Pid)
            )       
    }),
    session(Pid, Parent).
\end{verbatim} 

\noindent A suitable predicate API hiding the details of the protocol from the programmer can be written like so:

\begin{verbatim}
toplevel_call(Pid, Template, Goal) :-
    Pid ! '$call'(Template, Goal).

toplevel_next(Pid) :-
    Pid ! '$next'.

toplevel_stop(Pid) :-
    Pid ! '$stop'.
\end{verbatim} 

\noindent Here is a session that shows that these predicates work as intended:

\begin{verbatim}
?- toplevel(Pid, [
       monitor(true)
   ]).
Pid = 23891373.
?- toplevel_call($Pid, Who, human(Who)).
true.
?- flush.
Shell got success(23891373,socrates,true)
true.
?- toplevel_next($Pid).
true.
?- flush.
Shell got success(23891373,plato,true)
true.
?- toplevel_next($Pid).
true.
?- flush.
Shell got success(23891373,aristotle,false)
true.
?-
\end{verbatim}

\noindent Note that because of the outer recursive loop, our simple toplevel actor is again poised to run another query:

\begin{verbatim}
?- toplevel_call($Pid, X, q(X)).
true.
?- flush.
Shell got down(23891373,23891373,existence_error(proc,q/1))
true.
?-
\end{verbatim}

\noindent Here the actor crashed, and the reason is that the code for \texttt{toplevel/1-2} does not say anything about what should happen if an error is thrown, which is the case if the predicate called by the goal is not defined. Since the spawned process is monitored, however, the error message did eventually reach the mailbox of the spawning process anyway, in the form of a \texttt{down} message.

A similar problem is evident should we tell the actor to execute a query that does not terminate. Using \texttt{exit/2} works, but will once again kill the entire actor process.

Such problems can be fixed. As it turns out, however, and as will be demonstrated in Chapter~\ref{sec:prolog-agents}, Web Prolog programmers do not need to define their own toplevel predicates. Instead, they can use a set of built-in predicates in order to create and control more powerful Prolog toplevels, which are actors adhering to a much improved protocol. In the terminology introduced above, these are Erlang-style behaviors implementing a standardized toplevel protocol.

\section{Summary}\label{sec:discussion-revised}

This chapter introduced Web Prolog as a hybrid language that extends a substantial subset of ISO Prolog with Erlang-style constructs for programming concurrency. The core additions -- actors with private state, asynchronous message passing (\texttt{!/2}), process creation (\texttt{spawn/2-3}), and fault-handling primitives such as links, monitors, and exit signals. A program that never mentions these constructs behaves like a conventional Prolog program, while programs that do use them gain access to an Erlang-like world of concurrent processes.

We also argued that Erlang provides an unusually suitable point of departure for such an extension. Its actor model is mature, battle-tested, and close enough in syntax and spirit to Prolog that both languages can share patterns, idioms, and even didactic material. Many of the examples in this chapter are thin Prolog translations of familiar Erlang patterns. The combination of Prolog's backtracking search with Erlang-style concurrency is made workable by allowing receive to be semi-deterministic and by separating the nondeterminism of search from the determinism of message handling. This makes it possible to treat a Prolog toplevel as just another actor and to encapsulate interactive sessions as behaviors implemented on top of a common concurrency substrate.

Taken together, these choices define Web Prolog as a multi-paradigm language. Its logical expressivity remains that of ISO Prolog, but its operational repertoire is enlarged by an actor-based concurrency model that is both network-transparent and compatible with the rest of the Prolog Trinity. There are other possible ways of adding concurrency to Prolog, but the thesis of this chapter is simply that Erlang-style actors provide a coherent and practically manageable foundation on which the subsequent chapters can build.

%\subsection{The syntax and semantics of receive}
%
%\subsubsection{The syntax of receive}
%
%Once again, we show a priority queue example borrowed from Fred H\'ebert's textbook on Erlang \citep{Hebert:2013:LYE:2543986}. The purpose is to build a list of messages with those with a priority above 10 coming first:
%
%\begin{figure}
%\begin{minipage}{0.485\linewidth}
%\begin{tcolorbox}[
%  colback=white,
%  colframe=black,
%  arc=0pt,
%  outer arc=0pt,
%  boxrule=0.4pt,
%  left=2mm,right=2mm,top=1mm,bottom=1mm,
%  enhanced
%]
%\small
%\begin{verbatim}
%important() ->
%    receive
%        {P,M} when P > 10 ->
%            [M | important()]
%    after 0 ->
%        normal()
%    end.
%    
%\end{verbatim}
%\end{tcolorbox}
%\end{minipage}
%\hfill
%\begin{minipage}{0.485\linewidth}
%\begin{tcolorbox}[
%  colback=white,
%  colframe=black,
%  arc=0pt,
%  outer arc=0pt,
%  boxrule=0.4pt,
%  left=2mm,right=2mm,top=1mm,bottom=1mm,
%  enhanced
%]
%\small
%\begin{verbatim}
%important(Ms) :-
%    receive({
%        P-M if P > 10 ->
%            Ms = [M | MoreMs],
%            important(MoreMs)
%    },[ timeout(0),
%        on_timeout(normal(Ms))
%    ]). 
%\end{verbatim}
%\end{tcolorbox}
%\end{minipage}
%\end{figure}
%
%
%% \begin{verbatim}
%% important(Ms) :-
%%     receive({
%%         P-M if P > 10 ->
%%             Ms = [M|MoreMs],
%%             important(MoreMs)
%%     },[ timeout(0),
%%         on_timeout(normal(Ms))
%%     ]).
%% \end{verbatim}
%
%% \noindent For comparison, here is H\'ebert's version of \texttt{important/0}:
%
%% \begin{verbatim}
%% important() ->
%%     receive
%%         {P, Message} when P > 10 ->
%%             [Message | important()]
%%     after 0 ->
%%         normal()
%%     end.
%% \end{verbatim}
%
%\noindent Compared to the Web Prolog predicate, the Erlang function is more succinct. There are three reasons for this. First, Web Prolog has a relational syntax which does not allow nesting of calls while Erlang is a functional language where such nesting is the norm. 
%
%Secondly, while \texttt{op/3} allowed us to define the infix \texttt{if} operator, not much can be done about the more complex \texttt{receive...after...end} construct. Instead, Web Prolog defines a binary predicate expecting an ordered sequence of receive clauses wrapped in curly brackets in its first argument and the \texttt{after...} part as a list of options specifying the behaviour around time-outs in the second argument.
%
%Thirdly, while we could have taken advantage of the fact that an Erlang tuple such as \texttt{\{P,\,M\}} is valid syntax in Prolog too, where it is a somewhat more complex compound term of the form \texttt{\{\}((P,\,M))}, we chose not to. In Prolog it is always better to use less complex terms, taking up less memory. The ``pair operator'' (\texttt{-}) was a sensible choice here, but any valid Prolog term would do. %can be used in a pattern.
%
%Three fairly major syntactic differences in such a small program -- where at least the first two contribute to the program looking neater in Erlang than in Web Prolog -- may seem like a lot for a language that tries to appear as similar as possible to Erlang. However, note that the example was chosen exactly because it highlights many differences in a tiny program. Usually, the differences are less conspicuous.
%
%%All in all, we believe that most Erlangers will find the syntax of Web Prolog likeable and easy to work with. Fans of Elixir may be less enthusiastic, at least if the Prolog-ish syntax that Erlang inherited from Prolog is what drew them to Elixir with its Ruby-ish syntax instead. %(Of course, other factors such as tooling may have been at play here as well.)
%
%\subsubsection{The semantics of receive}
%
%Semantically, the receive operation must be specified precisely enough that implementers can build predictable selection and scheduling strategies.
%
%%\subsubsection{Intended semantics of \texttt{receive/1-2}.}
%
%Web Prolog's \texttt{receive} is defined in terms of \emph{selection}, \emph{consumption}, and \emph{execution}, with a clear boundary between tentative matching and irreversible mailbox update.
%
%\begin{enumerate}
%\item \textbf{Selection.}
%The runtime scans the mailbox and searches for a message $M$ that can be handled by some receive clause. For each candidate message, the clause head's pattern is tried by unification, and any associated guard is evaluated. Unifications performed while testing candidates are \emph{tentative} and are discarded unless a candidate is ultimately selected.
%
%\item \textbf{Consumption.}
%Once a particular message $M$ and clause have been chosen (that is, head unification succeeded and the guard, if present, succeeded), $M$ is removed from the mailbox. Note that a guard is tested \emph{before} message consumption. If head unification succeeds but the guard fails, the message is not consumed and the search continues.
%
%\item \textbf{Execution.}
%The clause body is executed under the bindings produced by head unification, together with any bindings or constraints produced by the guard.
%\end{enumerate}
%
%\noindent Guards are intended to be terminating, semideterministic goals (they may bind variables, but should not enumerate). A guard is evaluated as if enclosed in  \texttt{once/1}. If the guard throws an exception, it is treated as guard failure (the clause is not applicable, and the message is not consumed).
%
%---
%
%A call to \texttt{receive} selects and consumes at most one message. If the selected body succeeds, \texttt{receive} succeeds; if the selected body fails, \texttt{receive} fails, and the selected message remains consumed. Backtracking never revisits selection or consumption: no backtracking step can cause \texttt{receive} to re-select the same message or to select a different message. Any nondeterminism arises solely from the selected clause body.
%
%Thus, variable bindings are undone on failure as usual in Prolog, but mailbox state is not undone: message consumption is irreversible and is never rolled back by backtracking.
%
%Erlang's \texttt{receive} is also a mailbox-selection construct, but it differs from Web Prolog's \texttt{receive} in three important ways. First, Erlang clause heads use \emph{pattern matching} with \emph{single-assignment} variables, whereas Web Prolog clause heads use \emph{unification} with Prolog's logic variables (and optional guards expressed as Prolog goals). Second, Erlang's \texttt{receive} is operationally \emph{committing and deterministic}: once a message has been selected, it is consumed and exactly one clause body is executed without any notion of backtracking over alternative outcomes. In Web Prolog, selection and consumption are likewise single-shot, but the selected clause body may be nondeterministic, so backtracking can enumerate multiple solutions produced by that body while never revisiting message selection. Third, Erlang's selective receive scans the mailbox for the first message that matches any clause, and matching has no observable intermediate bindings beyond the chosen clause; Web Prolog may tentatively unify against candidates while searching, but all such bindings are discarded unless a message is selected and its guard succeeds. In short, both languages separate mailbox selection from subsequent computation, but Erlang enforces a single, deterministic continuation, whereas Web Prolog allows nondeterminism \emph{within} the chosen continuation while keeping mailbox effects irreversible.
%
%---
%
%Erlang's receive is an expression (a built-in control construct): when a clause matches, its body is evaluated and the receive expression yields the value of that body (in practice, the value of the last expression in the clause body). By contrast, Web Prolog's receive is a goal: it succeeds or fails, it does not itself yield a value, and any information is returned via variable bindings and subsequent goals rather than as an expression value.
%
%\subsubsection{Arbitrary goals as guards}
%
%Web Prolog allows receive patterns to be refined by guards written as Prolog goals. This is an expressiveness win for protocol programming: acceptance of a message can depend on semantic conditions grounded in local knowledge, private state, or policy predicates, not only on message shape.
%
%The cost is equally clear. Erlang restricts guards precisely to keep selection predictable and to prevent mailbox scanning from turning into arbitrary computation. In Web Prolog, programmers and implementers must instead rely on discipline and specification. Guards must be terminating and cheap relative to mailbox traffic, and they should be treated as part of a protocol surface rather than as a general computation hook.
%
%This is not merely a matter of style. If guard evaluation can be expensive or non-terminating, then the language loses a key claim: that the concurrency model remains simple enough to teach and reason about.
%
%\subsubsection{Semi-determinism and failure as control}
%
%A second design choice is that receive is treated as semi-deterministic: it either fails, or succeeds exactly once. This makes receive compose naturally with Prolog control. In particular, it allows failure-driven loops and restart patterns to be expressed without inventing a special control regime for message reception.
%
%The key semantic stance is that a failing clause body causes the receive call to fail. This preserves ordinary Prolog meaning: failure is a first-class control signal. It also aligns with the Trinity's later agent patterns, where toplevel-like sessions and protocol loops benefit from the ability to use failure intentionally and locally.
%
%\subsubsection{Fairness, starvation, and disciplined receive}
%
%Selective receive always raises the spectre of starvation: if a mailbox contains a large prefix of messages that do not match the currently preferred patterns, a naive selection strategy can delay progress indefinitely. This is not unique to Web Prolog; it is a known tension in Erlang as well. The difference is that Prolog's extra expressiveness can amplify the issue if guards are expensive or if reception is embedded in search.
%
%The mitigation is partly semantic and partly stylistic. Semantically, the receive operation must be specified precisely enough that implementers can build predictable selection and scheduling strategies. Stylistically, programmers should treat a receive specification as a protocol: clause order should reflect priorities, guards should be cheap, and long-running work should be separated from the mailbox scanning phase. Where necessary, an explicit dispatcher pattern can be used to decouple selection from processing. These are the same instincts that make Erlang code robust, but they become more important in a setting where guards can perform arbitrary computation.
%
%
%
%\subsection{Separating notification from lifetime coupling}\label{sec:microexamples}
%
%\noindent Web Prolog makes a sharp distinction between \emph{being told that another process terminated} (monitoring) and \emph{having one process's lifetime depend on another} (the default parent-to-child dependency of \texttt{spawn}). Erlang, by contrast, historically assigns both roles to one primitive: \emph{linking}. Since many readers will be less familiar with Erlang, we begin with a minimal explanation of how linking interacts with exit signals and the \texttt{trap\_exit} flag.
%
%In Erlang, a \emph{link} is a symmetric relationship between two processes. When one linked process terminates, an \emph{exit signal} with a \emph{reason} is delivered to the other. The default behaviour is deliberately fail-fast: if a process receives an exit signal with a reason other than \texttt{normal}, it terminates as well. This makes failures propagate automatically across linked components and is the basis for Erlang's supervision philosophy: rather than attempting to recover locally, a crashing component takes down the linked part of the subsystem, allowing a supervisor process to restart the whole group in a known state.
%
%The reason \texttt{normal} is special. A process that terminates normally still emits an exit signal to linked peers, but a \texttt{normal} signal is ignored by processes that do not explicitly opt in to receiving such events. This choice makes linking useful as a lifetime relationship without turning ordinary completion into noise: most processes can finish normally without generating messages or cascading terminations.
%
%A process can opt in to receiving exit signals as ordinary messages by setting the process flag \texttt{trap\_exit} to \texttt{true}. With trapping enabled, an incoming exit signal is converted into an \texttt{EXIT} message of the form \texttt{\{'EXIT',\,Pid,\,Reason\}} and delivered to the process mailbox. Importantly, trapping changes how abnormal exits from linked peers are handled: instead of being terminated by the signal, the process is notified and can decide what to do. (The one exception is \texttt{kill}, which remains untrappable in Erlang, but that detail is not needed for the examples below.) Thus, in Erlang, \texttt{trap\_exit} is the mechanism that allows a process to \emph{observe} linked failures without automatically sharing their fate.
%
%With this background, we can now compare Erlang's ``link plus optional trapping'' pattern with Web Prolog's explicit separation into monitors (notification) and default one-way spawn dependencies (lifetime coupling).
%
%Seen from the perspective of this comparison, Erlang's mechanism is undeniably subtle. The meaning of a \emph{link} depends not only on the direction in which an exit happens, but also on the exit \emph{reason} (with \texttt{normal} treated specially) and on a per-process mode bit (\texttt{trap\_exit}) that changes whether an incoming exit is lethal or becomes a mailbox message. In other words, linking simultaneously serves as a tool for lifetime coupling, a tool for failure propagation, and -- conditionally -- a tool for notification, and the programmer must keep track of several interacting exceptions. This is not a criticism of Erlang: the design reflects a long history of hard-won engineering choices in a language whose primary mission is building robust systems.
%
%For our purposes, however, the conditionality is precisely what motivates a different approach. \cite{conf/erlang/SvenssonFE10} proposed a semantics in which the rules are made more uniform, with fewer special cases and fewer feature interactions. That work gave us both the inspiration and the courage to adopt a simpler model in Web Prolog: monitors are always just notification, and default spawn dependencies are always just one-way lifetime coupling. Throughout this book we have deliberately stayed close to that ``less buts and ifs'' design, not because Erlang is wrong, but because a small, orthogonal core is easier to explain, reason about, and reimplement across a federated Web setting.
%
%\subsubsection{Example A: abnormal child termination}\label{sec:microexample-abnormal-child}
%
%\begin{figure}[h]
%\begin{minipage}{0.485\linewidth}
%\begin{tcolorbox}[
%  colback=white,colframe=black,arc=0pt,outer arc=0pt,boxrule=0.4pt,
%  left=2mm,right=2mm,top=1mm,bottom=1mm,enhanced
%]
%\small
%\begin{verbatim}
%p() ->
%  process_flag(trap_exit,true),
%  P = spawn_link(fun c/0),
%  receive
%    {'EXIT', P, W} ->
%      format("~p ~p~n",[P,W])
%  end,
%  format("parent continues~n").
%
%c() ->
%  exit(error).
%\end{verbatim}
%\end{tcolorbox}
%\end{minipage}
%\hfill
%\begin{minipage}{0.485\linewidth}
%\begin{tcolorbox}[
%  colback=white,colframe=black,arc=0pt,outer arc=0pt,boxrule=0.4pt,
%  left=2mm,right=2mm,top=1mm,bottom=1mm,enhanced
%]
%\small
%\begin{verbatim}
%p :-
%  spawn(c, P),
%  monitor(P, R),
%  receive({
%    down(R, P, W) ->
%      format("~p ~p~n",[P,W])
%  }),
%  format("parent continues~n").
%
%c :-
%  exit(error).
%\end{verbatim}
%\end{tcolorbox}
%\end{minipage}
%\end{figure}
%
%\noindent In the Erlang example, \texttt{spawn\_link} creates a \emph{bidirectional} link: when \texttt{c} dies abnormally, an exit signal is sent to \texttt{p}. Unless \texttt{p} traps exits it would be terminated as well; with \texttt{trap\_exit=true} the signal becomes an \texttt{EXIT} message that can be printed, after which \texttt{p} continues. In the corresponding Web Prolog example, \texttt{p} does not need an analogue of \texttt{trap\_exit} to survive, because termination is observed via a monitor (\texttt{down/3}) rather than delivered as a lethal signal. This example isolates the first key difference: Erlang links combine notification with potential failure propagation, while Web Prolog monitoring provides notification without upward propagation.
%
%\subsubsection{Example B: normal parent termination}\label{sec:microexample-normal-parent}
%
%\begin{figure}[h]
%\begin{minipage}{0.485\linewidth}
%\begin{tcolorbox}[
%  colback=white,colframe=black,arc=0pt,outer arc=0pt,boxrule=0.4pt,
%  left=2mm,right=2mm,top=1mm,bottom=1mm,enhanced
%]
%\small
%\begin{verbatim}
%p() ->
%  spawn_link(fun c/0),
%  ok. % p terminates normally
%
%c() ->
%  timer:sleep(1000),
%  format("child still alive~n").
%\end{verbatim}
%\end{tcolorbox}
%\end{minipage}
%\hfill
%\begin{minipage}{0.485\linewidth}
%\begin{tcolorbox}[
%  colback=white,colframe=black,arc=0pt,outer arc=0pt,boxrule=0.4pt,
%  left=2mm,right=2mm,top=1mm,bottom=1mm,enhanced
%]
%\small
%\begin{verbatim}
%p :-
%  spawn(c),
%  true. % p terminates normally
%
%c :-
%  sleep(1),
%  format("child still alive~n").
%\end{verbatim}
%\end{tcolorbox}
%\end{minipage}
%\end{figure}
%
%\noindent In Erlang, when \texttt{p} terminates normally, the \texttt{normal} exit signal sent to a linked process is ignored by a non-trapping child, so \texttt{c} keeps running and prints after the delay. In Web Prolog, the default \texttt{spawn} dependency is \emph{one-way}: the child depends on the parent. Therefore, when \texttt{p} terminates with reason \texttt{true} (i.e. normally), \texttt{c} is terminated as well and the delayed print never occurs.\footnote{The delays are not part of the concurrency mechanism being compared. Its only purpose is to make the outcome observable. If the child attempts to print immediately, it may happen to run before the parent has terminated, and then both versions would appear to ``work'' even though their termination semantics differ.} This isolates the second key difference: Erlang's special treatment of \texttt{normal} allows linked processes to outlive each other, whereas Web Prolog's default parent-to-child dependency does not.
%
%
%\subsubsection{Example C: Child termination in Erlang and Web Prolog}\label{sec:microexample-normal-child}
%
%\begin{figure}[h]
%\begin{minipage}{0.485\linewidth}
%\begin{tcolorbox}[
%  colback=white,colframe=black,arc=0pt,outer arc=0pt,boxrule=0.4pt,
%  left=2mm,right=2mm,top=1mm,bottom=1mm,enhanced
%]
%\small
%\begin{verbatim}
%p() ->
%  P = spawn_link(fun c/0),
%  receive
%    {'EXIT', P, W} ->
%      format("~p ~p~n",[P,W])
%  after 200 ->
%      format("no EXIT message~n")
%  end,
%  format("parent continues~n").
%
%c() -> 
%  ok. % normal termination
%\end{verbatim}
%\end{tcolorbox}
%\end{minipage}
%\hfill
%\begin{minipage}{0.485\linewidth}
%\begin{tcolorbox}[
%  colback=white,colframe=black,arc=0pt,outer arc=0pt,boxrule=0.4pt,
%  left=2mm,right=2mm,top=1mm,bottom=1mm,enhanced
%]
%\small
%\begin{verbatim}
%p :-
%  spawn(c, P),
%  monitor(P, R),
%  receive({
%    down(R, P, W) ->
%      format("~p ~p~n",[P,W])
%  }),
%  format("parent continues~n").
%
%c :-
%  true. % normal termination
%
%\end{verbatim}
%\end{tcolorbox}
%\end{minipage}
%\end{figure}
%
%\noindent This example isolates a small but practically important distinction. In Erlang, linking does not guarantee a message-level notification of termination: if \texttt{p} does not set \texttt{trap\_exit=true} and \texttt{c} terminates normally, then the resulting \texttt{normal} exit signal is ignored and no \texttt{EXIT} message is delivered to \texttt{p}. (The \texttt{after} clause makes the absence of notification visible.) In Web Prolog, notification is the explicit purpose of monitoring: the monitor delivers \texttt{down(R,\,P,\,true)} even for normal termination, and \texttt{p} can react without having to enable any special ``trap'' mode.
%
%\subsubsection*{Summary: the invariant of one-way dependency}
%
%Examples A--C illustrate the guiding separation in Web Prolog: monitoring is solely a mechanism for termination notification, whereas spawn establishes the default lifetime relation. In Erlang, by contrast, linking historically bundles several roles into one primitive. A link is symmetric; whether it yields a mailbox-level notification depends on \texttt{trap\_exit}, and whether it propagates failure depends both on direction and on the exit reason, with \texttt{normal} treated specially. Web Prolog deliberately avoids this conditionality: monitors always notify, and the default spawn dependency is always one-way lifetime coupling. 
%
%The resulting invariant is simple and reason-agnostic: when a parent process ceases to exist, its children are pruned regardless of whether the parent exited normally (with reason \texttt{true}), failed (with reason \texttt{true}), reported an error as its reason, or was killed. This is an architectural shift rather than a surface convenience. Erlang links are often used to create meshes of peers where failure propagates but ordinary completion does not, while Web Prolog treats the process hierarchy as a scope for concurrent activity. A spawn is therefore closer to a nested block in conventional programming: when the block is exited, the resources created inside it are reclaimed. This shift aligns Web Prolog with modern structured concurrency principles, where the lifetime of concurrent tasks is bound to the blocks that created them.\footnote{See e.g. \url{https://en.wikipedia.org/wiki/Structured\_concurrency}} 
%
%% Table~\ref{table:structured-concurrency} summarizes the contrast. 
%
%% \begin{table}[h]
%% \centering
%% \small
%% \begin{tabularx}{\textwidth}{@{\hspace{20pt}}l@{\hspace{20pt}}X@{\hspace{20pt}}X@{}}
%% \hline
%% \textbf{Parent Exit Reason} & \textbf{Erlang Child (Linked)} & \textbf{Web Prolog Child} \\ \hline
%% \texttt{normal / true}      & Continues running              & Terminated                \\
%% \texttt{false}              & Not applicable                 & Terminated                \\
%% \texttt{error(Reason)}      & Terminated (if not trapping)   & Terminated                \\
%% \texttt{kill}               & Terminated (untrappable)       & Terminated                \\ \hline
%% \end{tabularx}
%% \caption{Comparison of child survival based on parent exit reason.}
%% \label{table:structured-concurrency}
%% \end{table}
%
%The key consequence is that Web Prolog makes resource reasoning straightforward: no background activity outlives the process that initiated it unless the programmer explicitly chooses to decouple it (for example by delegating the spawn to a long-lived ancestor or by using a detached flag, if the implementation provides one). 
%
%
%
%\subsubsection{Achieving peer independence}\label{sec:peer-independence}
%
%While the default behavior of Web Prolog is hierarchical, there are cases where a process must outlive its creator -- for instance, a logger or a background service that should persist even if the request-handling process that initiated it finishes. To achieve this 'peer independence', the activity must be detached from the local scope.
%
%In Web Prolog, this is typically handled by delegating the spawn to a long-lived ancestor or using a 'detached' flag (if the implementation provides one). This effectively moves the child up the process tree, as illustrated below:
%
%\begin{figure}[h]
%\centering
%\small
%\begin{tabular}{p{0.45\linewidth} | p{0.45\linewidth}}
%\textbf{Standard Hierarchy} & \textbf{Detached / Peer Behavior} \\ \hline
%Parent spawns Child. Child is a leaf of the Parent. If Parent dies for any reason, Child is pruned. & Parent requests a 'System' process to spawn Child. Child is now a sibling to the Parent. Parent can die while Child continues under the System's protection. \\
%\end{tabular}
%\end{figure}
%
%This architectural choice forces the programmer to be explicit about process longevity. In Erlang, longevity is the default (unless a link is created and a failure occurs); in Web Prolog, cleanup is the default, and longevity must be requested. This shift aligns Web Prolog with modern 'structured concurrency' principles, where the lifecycle of concurrent tasks is strictly bound to the blocks that created them.

% ===

% This concludes the initial presentation of the concurrency-oriented features supported by the Web Prolog profile, and a summary and discussion is in order. The argument has been made that extending a large subset of ISO Prolog with Erlang-style constructs for programming concurrency creates an excellent multi-paradigm programming language. 

% We have toyed with the idea of viewing Web Prolog as a dialect of Erlang with logic variables, backtracking search and a built-in database of facts and rules -- an idea not too farfetched. Indeed, if it was not for the fact that Web Prolog resembles Prolog a lot more than it resembles Erlang, we might have named it ``Web Erlang.'' But we did not, and the reason is that, with the exception of the Erlang-style constructs, Web Prolog stays really close to traditional Prolog, indeed so close that the vast majority of example programs in Prolog text books should run without modification. (When they do not run, it is usually because the book in question does not follow the ISO standard.) 

% We have explained that one reason we choose to go with Erlang-style concurrency rather than something different is that Erlang is arguably the most mature concurrency-oriented language in existence. We have praised the maturity of Erlang, but maturity can also carry a risk. Design choices that were once pragmatic may harden into permanent features -- even when they no longer serve the language well.

% Therefore, the design does not slavishly follow Erlang. We have chosen to simplify Web Prolog's Erlang-style features in certain ways, especially as regards the machinery around links and monitors. But we have also chosen to extend it in some ways, with features of the receive primitive that Erlang does not support. We will discuss these choices in separate sections.

% In this section we summarize our design choices and why we deviate from Erlang in a few places.

% ---

% justify going beyond and abstracting over the ISO Prolog threads

% ---

% So the semantics of concurrency-oriented primitives, many design patterns, the thinking in terms of protocols, and the use of programming abstractions (behaviors) were deemed worth adapting to Web Prolog. 

% Also, for reasons we have stated, Erlang happens to be syntactically very similar to Prolog. We try to make sure that Web Prolog stays as similar to Erlang as possible, not only syntactically, but also when it comes to semantics and pragmatics. Such similarities means that Erlang-style concurrent programming in Web Prolog can be learned from the many text books that has been published by members of the Erlang community. 

% ---

% As we saw in this chapter, just like in Erlang, actors are created by calling \texttt{spawn/2} and they communicate by sending and receiving messages. Messages are sent by using the send operator \texttt{!/2} and are selected from the mailbox of the recipient by calling \texttt{receive/1}. The spawn operation as well as the receive can be configured by means of options, so there is a \texttt{spawn/3} and a \texttt{receive/2} as well, where the optional last argument are lists of options, some of which were described in this chapter. 

% By design, and just like in Erlang, the communication between two concurrently running Prolog actors is asynchronous. It means that one actor may send the other actor a message and can then just carry on with other tasks without first waiting for a response to arrive. Also, just like in Erlang, sending is by default ``fire and forget'' in the sense that the sending actor will often neither know nor care whether the message sent has reached the client. (In \textit{Programming Erlang}, Armstrong refers to this as message passing with ``send and pray semantics.'') However, if a response \textit{is} necessary for the sake of smooth client-actor cooperation, it can easily be arranged. On top of asynchronous communication it is usually fairly straightforward to build an actor which communicates in a synchronous manner, i.e. an actor which after having sent a message is \textit{programmed} to wait idly for a response to arrive. 

% It all means that porting Erlang programs to Web Prolog should be fairly straightforward. In fact, many of the programs used as examples in this chapter have been ported from examples appearing in Erlang text books. (A larger program ported from Erlang into Web Prolog can be found in Appendix~\ref{sec:big-example}.)

% ---

% Treating \texttt{!/2} and \texttt{exit/2} as no-ops when the pid in the first argument points to a non-existent actor process.

% \subsection{We are not the first to be inspired by Erlang}\label{sec:not-the-first}

% ---

% Not the first too base our stuff on erlang
% Examples - ask Gemini!
% But we did it better
% ISO Threads
% Deviations from Erlang.

% \subsection{A hybrid identity: The best of both worlds}\label{sec:hybrid-identity-wp}

% Web Prolog's identity is intentionally hybrid. It can be viewed from two perspectives:

% \begin{itemize}
%     \item \textbf{For the Prolog programmer}: It is ``good ol' Prolog'' augmented with a principled, battle-tested model for concurrency and distribution. It retains the core features of ISO Prolog, ensuring that the vast majority of existing code and textbook examples run without modification. The Erlang-style features are an \textit{extension} to its procedural capabilities, not a replacement of its logical core.
%     \item \textbf{For the Erlang programmer}: It is a dialect of Erlang where the sequential, functional core has been replaced with a relational, logic programming engine complete with backtracking, unification, and a built-in database. This substitution, as Armstrong himself noted was theoretically sound, opens the door to solving problems in knowledge representation and reasoning within the actor model.
% \end{itemize}

% \noindent This fusion was achieved by adhering to a key design principle: syntactic and semantic consistency. To a large extent, Web Prolog adopts Erlang's naming conventions and operational semantics for its concurrency primitives, allowing developers to leverage the wealth of existing Erlang literature and design patterns. However, our design deviates from Erlang in a couple of areas.

% \subsection{A more powerful receive construct}\label{sec:powerful-receive}

% \begin{quotation}
% I would argue that it is precisely the 'receive' construct in Erlang that makes Erlang such a joy to use.\footnote{\url{https://groups.google.com/forum/\#!msg/erlang-programming/gjU-HCoq7dk/Mx\_Af0iQ5P0J}}
% \vspace{-0.3cm}
% \flushright \textit{Richard O'Keefe}
% \end{quotation}

% \noindent We are very happy with our adaptation of Erlang's receive operator. It is, just like in Erlang, a joy to use. Also, adding to this joy, the receive operation in Web Prolog has two powerful features that Erlang lacks: we allow arbitrary goals in guards, and we allow a call to \texttt{receive/1-2} to fail.

% \subsubsection{Arbitrary goals as guards}\label{sec:arbitrary-goals-as-guards}

% [A COPY OF THIS IS PLACED IN THE SM secton!]:

% Erlang enforces purity and efficiency by only allowing a restricted set of primitives in guards, and completely disallows calling user-defined functions. For the use cases the people behind Erlang had, it probably made sense to impose such restrictions. 

% In Web Prolog, however, although only the first solution will be searched for, \textit{any} goal may be used as a guard and values of variables bound by it are available in the body. This makes the receive construct more powerful than in Erlang, but it also means that the programmer is made responsible for keeping guards as simple and efficient as possible and to avoid side effects. In our opinion, and as exemplified by the state-machine code in Section~\ref{sec:implementing-state-machines}, enabling the programmer to condition the matching of a receive clause on the content of the entire Prolog database makes it worth it. 

% Most Erlang developers accept that guards serve a specific role, and they understand the trade-offs. But some find it limiting and express desire for more expressive, reusable guard logic. So there is broad interest and research into relaxing the restriction for \textit{pure} functions.

% ---

% The change of operator names from \texttt{when} to \texttt{if} suggests the change of semantics. (This is not the only reason why we prefer \texttt{if}.) 

% \subsubsection{Semi-deterministic behavior}\label{sec:semi-deterministic-behavior}

% A feature that the receive predicate in Web Prolog has, but Erlang lacks, is that it is a \textit{semi-deterministic} predicate, i.e. it either fails, or succeeds exactly once. This is a key property of \texttt{receive/1-2} which ensures that backtracking can be handled in an elegant way. Of course, this would not even make sense in Erlang, as it is deterministic by design and does not do backtracking.

% ---

% Section~\ref{sec:answers-through-backtracking} and \ref{sec:simple-toplevel}

% ---

% As will be shown in Chapter~\ref{sec:prolog-agents}, it is also what allows the introduction of first-class Prolog toplevels into the language.

% ---

% For an Erlang programmer this particular use of \texttt{receive/1-2} may come as a surprise. After all, the Prolog concepts of \textit{failure} and \textit{backtracking} and the use of failure to force backtracking are foreign to Erlang. Prolog programmers may recognize a behavior due to the fact that \texttt{receive/1-2} is a \textit{semi-deterministic} predicate, i.e. a predicate that either fails, or succeeds exactly once. We may want to compare \texttt{receive/1} with \texttt{read/1} in ISO Prolog, which is also semi-deterministic, and also serve the purpose of receiving data from the environment. But while \texttt{read/1} fails if the pattern (consisting of any term) in the argument does not unify with the term that is read, the only way \texttt{receive/1-2} will fail is if the goal in the \textit{body} of one of its receive clauses fails, or if timeout happens and the goal passed in the \texttt{on\_timeout} option fails. We would suggest that this is the only reasonable semantics for \texttt{receive/1-2}. What else could it possibly mean if the goal in the body of a receive clause fails? Raising an error?

% We would argue that letting the whole call to \texttt{receive/1-2} fail gives us the only reasonable semantics. 

% Before the failure occurs, operations with side effects can be performed -- operations such as sending a message (as we did in sections \ref{sec:answers-through-backtracking} and \ref{sec:simple-toplevel}), perform I/O, or update the private dynamic database.

% \subsection{Our simple model for links and monitors}\label{sec:simpler-links}

% \todo{Talk to Hans Svensson or Koen Claessen about this section and next!}
% Compared to Erlang, Web Prolog offers a simpler model for links and monitors. Where Erlang retains historical baggage in the form of bidirectional links and the \texttt{trap\_exit} mechanism, Web Prolog avoids these complications and relies entirely on unidirectional links combined with monitoring. This not only simplifies the programmer's mental model, but also reduces the potential for subtle errors.

% To aid in a comparison between Web Prolog and Erlang, consider the two programs below -- a Web Prolog program to the left and an Erlang program to the right:

% {\small
% \begin{verbatim}
% demo :-                             demo() ->                       
%   spawn(child, P),                    process_flag(trap_exit, true),
%   monitor(P, R),                      P = spawn_link(fun() ->           
%   receive({                                            child()          
%     down(R, P, W) ->                                 end),              
%       format("~p ~p~n",[P,W])         receive                       
%   }).                                   {'EXIT', P, W} ->           
%                                           io:format("~p ~p~n",[P,W])
% child :-                              end.                          
%   spawn(receive({a -> true})),                        
%   exit(crashed).                    child() ->                      
%                                       spawn_link(                   
%                                         fun() -> grandchild() end 
%                                           receive a -> ok end
%                                         end),                       
%                                       exit(crashed).                   
% \end{verbatim}
% }

% \noindent The running of these programs involves three actors each: a parent (i.e. the shell), a child, and a grandchild. Here is what happens when \texttt{demo/0} is called in the Web Prolog shell and in the Erlang shell, respectively:

% {\small
% \begin{verbatim}
%   ?- demo.                             2> demo:demo().    
%   51787912 crashed                     <0.175.0> crashed  
%   true.                                ok                 
%   ?-                                   3>                 
% \end{verbatim}
% }

% \noindent The child kills itself by calling \texttt{exit(crashed)}, and since the child links to the grandchild, the grandchild dies too. The life of the parent is not affected -- in the Web Prolog program because links are unidirectional, and in Erlang because the exit signal is trapped and converted into an exit message.

% What happens if the child and the grandchild are not linked? In Web Prolog this is accomplished by changing the line \texttt{spawn(spawn(receive(\{\}))} into \texttt{spawn(spawn(receive(\{\})),\,\_,\,[link(false)])}. In Erlang we just use \texttt{spawn/1} instead of \texttt{spawn\_link/1}. In both scenarios the effect is that the grandchild process becomes an \textit{orphan}.

% What happens if the grandchild is crashed? In Web Prolog, if we replace \texttt{receive(\{\})} with \texttt{exit(crashed)} the child terminates normally. In Erlang, if we replace the forever-blocking receive with \texttt{exit(crashed)} the child crashes abnormally. But if the \texttt{exit} signal is trapped here too, it terminates normally.

% ---

% In Erlang, the trapping of the \texttt{exit} signal and the conversion into an \texttt{exit} message is crucial. If it is removed, then no \texttt{exit} message will be produced, and what is more, the \texttt{exit} \textit{signal} would reach the shell and kill it.

% In Web Prolog, the death of a spawned actor process will only generate at most one kind of built-in message -- the \texttt{down} message -- and only if the actor is monitored.

% ---

% This is intentionally asymmetric, and links are no longer ``mutual death pacts'' as in Erlang. Instead, the link defines who depends on whom.

% ---

% perhaps use receive a -> ok end in both programs.

% The overlap couldn't be illustrated more clearly. 

% Web Prolog does not have exit messages, and so always have to depend on monitors.

% ---

% As the code examples in this chapter clearly show, the building of supervision hierarchies appears to be possible.

% ===

% \subsubsection{Killing the shell}

% What would happen if we remove the call to \texttt{process\_flag(trap\_exit,\,true)} in the Erlang program?

% ---

% In our Erlang \texttt{demo/0}, the shell sets 
% \texttt{process\_flag(trap\_exit,true)}. This ensures that the child's abnormal 
% termination is delivered as a message instead of killing the shell. If we remove this trap, however, the behavior changes dramatically: the shell is linked bidirectionally with the child, so when the child calls \texttt{exit(crashed)}, the shell is immediately taken down by the exit signal. 
% The waiting \texttt{receive} never runs, and nothing is printed. 

% ---

% Note that in Erlang, calling \texttt{spawn\_link/1} causes the shell to crash, while calling \texttt{spawn/3} with the \texttt{link} option set to \texttt{true} in Web Prolog does not lead to a crash. The behavior of Erlang is not much liked by Svensson et al. in \cite{conf/erlang/SvenssonFE10}, and they ``find it a bit strange that a  process can affect the behavior of another process behind its back by linking to it.''

% \subsubsection{Overlapping mechanisms}

% According to the authors, ``the overlap between monitors and links'' in current Erlang complicates the semantics without real benefit.

% The most striking example is the overlap  between monitors and links, where there is very little practical difference between a trapped link-message and a monitor-message. 

% Adding the following calls just after the call to \texttt{spawn\_link/1} would have produced the same printed output twice:

% {\small
% \begin{verbatim}
% erlang:monitor(P, R),
% receive
%   {'DOWN', R, process, P, W} ->
%     io:format("~p ~p~n",[P,W])
% end,  
% \end{verbatim}
% }

% \subsubsection{Scrap}

% \noindent Placed side by side, the two \texttt{demo/0} programs let us show that the 
% \emph{observable} behavior can be identical while the failure semantics and 
% the design intent differ. In Erlang, the parent links bidirectionally to its 
% child and sets \texttt{trap\_exit}, so the child's abnormal termination arrives 
% as a message while the link would otherwise be fatal; the grandchild dies 
% because it is linked to the crashing child. In Web Prolog, the parent instead 
% combines a \emph{unidirectional} link (for fate-sharing in one direction) with a 
% \emph{monitor} (for termination notification), and the child links downward to 
% ensure its grandchild is taken out on crash. Thus, the parent still prints the 
% child's pid and the reason, and the forever-blocked grandchild is still cleaned 
% up. But the decomposition into \emph{link for liveness coupling} and 
% \emph{monitor for information} is deliberate: it removes the ambiguity of 
% Erlang's trapped links and makes the causal chain explicit in the program text. 
% This is exactly the contrast our two snippets bring out succinctly. 

% Their semantics also treats all side effects (linking, monitoring, 
% spawning) as asynchronous and ``everything as distributed,'' which matches Web 
% Prolog's web-native setting and cleanly supports the monitor-plus-link pattern 
% used in our Web Prolog demo. In their presentation, a link is simply a one-way 
% obligation (if $A$ links to $B$ and $A$ dies, $B$ must die), while a monitor is 
% a one-way observation channel that delivers a \texttt{down} notice when $B$ 
% terminates. Our Web Prolog demo is a textbook instance of this: the child's 
% crash triggers both the downward fate-sharing to the grandchild and an 
% informational \texttt{down} to the parent, with no need for 
% \texttt{trap\_exit}. The paper explicitly motivates this move as a simplification 
% that remains adequate for supervisory structures once those structures are 
% phrased in terms of unidirectional links plus monitors -- precisely the discipline our 
% example exhibits. In short: the Web Prolog approach will work, and it aligns 
% with a principled, cleaner semantics advocated by Svensson et al. for a future 
% Erlang-like language such as Web Prolog.

% ---

% In Web Prolog, the corresponding \texttt{demo/0} is unaffected by such a change. The shell relies on a \emph{monitor} for termination notification, and the link is unidirectional (from shell to child). When the child exits, the shell is not dragged down, but simply receives a \texttt{down(Ref,\,ChildPid,\,crashed)} message from the monitor. In other words, the Erlang shell must trap to survive, while the Web Prolog shell is safe by construction.

% ---

% In the semantics for current Erlang, tightly coupled to legacy implementation details from  older versions of the Erlang runtime, processes are able to trap \texttt{exit}-signals, and treat them as ordinary messages, thus avoiding termination upon  receiving an \texttt{exit}-signal.   

% The author instead propose to only have unidirectional links, and to remove the whole \texttt{exit}-trapping functionality since a monitor serves the same purpose.  

% ---

% \noindent Links in Erlang are bidirectional, exit signals can be trapped and converted into messages, and there are programmatic ways to remove links.\footnote{For an excellent introduction to how links and monitors work in Erlang, see for example \url{https://learnyousomeerlang.com/errors-and-processes}} Compared to what Erlang provides, the machinery around links in Web Prolog is much simpler. Links in Web Prolog are unidirectional; parents link to their children, making their lives depend on the life of the parent, but not \textit{vice versa}. Trapping of exit signals is not necessary, and there is no way to remove a link. 

% ---

% The paper \emph{A Unified Semantics for Future Erlang} proposes a revised semantic model for Erlang, designed to be cleaner and conceptually simpler than the current formal semantics, which are tightly coupled to legacy implementation details from older versions of the Erlang runtime. 

% According to the authors, the existing semantics are overly complicated, primarily due to distinctions between local and distributed execution, and support for legacy constructs such as bidirectional links and the \texttt{trap\_exit} mechanism.

% ---

% The shell process serves as the grandparent to the child.

% ---

% The authors argue that they provide a more realistic and scalable foundation for future Erlang implementations, especially in multi-core and highly parallel environments.

% ---

% As argued by Svensson et al. in \cite{conf/erlang/SvenssonFE10}, unidirectional links simplify things and do not harm expressivity. They also argue that the exit trapping functionality of Erlang is superfluous since a monitor serves the same purpose. As our code examples show, the building of supervision hierarchies appears to be possible.

% Our streamlined model provides the necessary foundation for building supervision hierarchies without introducing the complexities of bidirectional dependencies, aligning with the principle of ``let it crash'' in a web context where parent processes (e.g., a client session) naturally dictate the lifecycle of their children.

% \subsection{Translating Erlang into Web Prolog and vice versa?}

% We have suggested that if we can handle the translation of a functional program into Prolog's relational model, and if we use the Erlang-style primitives for spawning and messaging that Web Prolog supports, then porting just about any Erlang program into a Web Prolog program that does the same thing should be feasible. Appendix ? contains a slightly bigger example.

% The value would be that the OTP can be portable. statem, client server, etc. Will be difficult though. 

% Is translatability a goal we should entertain, then it is good to keep in mind that every time we do something different than Erlang (and if Erlang also does it) we make porting more difficult. With the discrepancies that do exist between Erlang and Prolog, is hard enough as it is.

% No, because porting the OTP would be a one-shot thing.

% A DCG for parsing Erlang to a level that would allow us to translate Erlang source code into Web Prolog. 

% ---

% Most uses of receive in Erlang can be mechanically translated into uses of receive in Web Prolog that will behave in the same way as in Erlang. The other way around is not always possible. 

% ---

% Another reason the translation of a Web Prolog program into Erlang may in general not be possible is that the receive construct in Web Prolog is a \textit{semi-deterministic} predicate.

% \newpage

% \subsection{Comparison with Erlang}

% \begin{table}[h]
% \centering
% \small
% \setlength{\tabcolsep}{3pt}
% \renewcommand{\arraystretch}{1.2}

% \begin{tabularx}{12cm}{%
%   >{\raggedright\arraybackslash}p{2.9cm}
%   >{\raggedright\arraybackslash}X
%   >{\raggedright\arraybackslash}X
% }
% \hline
% \textbf{Feature} &
% \textbf{Web Prolog} &
% \textbf{Erlang}
% \\ \hline

% Isolation &
% Actor isolation, but ``inside'' each actor you still have full Prolog (including nondeterminism and shared \emph{logical} variables \emph{within} a single actor). &
% Actor isolation, with a strictly deterministic single-threaded process model; no shared variables even within a process.
% \\ \hline

% Network-transparency &
% Same programming model intended for local and remote actors; tends to emphasise web-facing interaction patterns  as first-class use cases. &
% Distribution is deeply integrated (pids, registered names, links/monitors across nodes), but is a BEAM/cluster feature rather than ``web-native'' by default.
% \\ \hline

% Asynchronous send &
% Yes; message terms are Prolog terms, and can naturally carry structured data (including nested terms) and correlate with goals and templates. &
% Yes; messages are Erlang terms; pattern matching is strong but does not include logic-variable unification.
% \\ \hline

% Send never throws &
% Intended to be non-throwing in the common pid case; errors are typically shifted to observation mechanisms (monitor/link) rather than synchronous send failure. &
% In the normal case, \texttt{!/2} does not raise when the receiver is gone; send failures are not reported synchronously (unless you do something explicitly error-prone, such as sending to an invalid destination).
% \\ \hline

% Selective receive &
% Selective receive with unification-based matching; can feel ``more declarative'' because patterns may contain logic variables that become bound by the match. &
% Selective receive with pattern matching (no unification); patterns test structure and bind variables once, but do not solve constraints.
% \\ \hline

% Guarded reception &
% Guards can be \emph{general Prolog goals} (potentially calling user predicates), which is very expressive but also makes receive selection depend on arbitrary computation. &
% Guards are deliberately restricted (side-effect–controlled, terminating tests), which improves predictability and compilation opportunities.
% \\ \hline

% Linking &
% Supported; often discussed as policy-driven (for example, whether exits propagate, and whether trapping is default or optional). &
% Core primitive; bidirectional linking plus exit signals; \texttt{process\_flag(trap\_exit,true)} controls whether exits kill you or arrive as messages.
% \\ \hline

% Monitoring &
% Spawn-time monitoring plus post-hoc \texttt{monitor/2}; the post-hoc case is useful but not atomic with respect to the target's lifetime. &
% \texttt{spawn\_monitor} gives atomic spawn+monitor; plain \texttt{monitor} is post-hoc and likewise non-atomic, but is widely used with the same caveat.
% \\ \hline

% \end{tabularx}
% \end{table}


% \subsection{Comparison with the ISO Prolog Threads proposal}\label{sec:comparison-iso-threads-wp}

% To better understand the design choices behind Web Prolog, it is instructive to compare its concurrency primitives with the ISO Prolog Threads working draft. Even when we restrict attention to a single machine, the two approaches reflect different default assumptions about state, communication, and failure handling. The ISO draft adopts a low-level, \textsc{posix}-style model based on threads, shared facilities, and explicit synchronization, while Web Prolog adopts an actor model in which isolated processes interact exclusively by asynchronous message passing.

% \begin{table}[ht]
% \centering
% \small
% \setlength{\tabcolsep}{3pt}
% \renewcommand{\arraystretch}{1.2}
% \begin{tabularx}{\textwidth}{l X X}
% \hline
% \textbf{Feature} & \textbf{Web Prolog (Actors)} & \textbf{ISO Threads (Shared-Memory)} \\ \hline
% Isolation & Full isolation: private mailboxes and encapsulated state. & Partial: shared dynamic code and global resources. \\
% Network & Designed for local and remote transparency. & Local only; distribution requires external protocols. \\
% Send Semantics & Non-throwing 'fire-and-forget' to Pids. & Often throws error on non-existent destinations. \\
% Receive Logic & Selective receive; unmatched messages are deferred. & FIFO; messages are consumed upon unification. \\
% Guards & Supported: arbitrary goals can guard reception. & No native guarded receive; requires manual dispatch. \\
% Fault Tolerance & Simplified Erlang-style linking and monitoring. & Manual join patterns and signal handling. \\ \hline
% \end{tabularx}
% \caption{Technical Comparison: Web Prolog vs. ISO Threads}
% \end{table}

% \subsubsection{Programming model and state}

% A critical distinction lying at the heart of this comparison is the concept of \emph{isolation}. In the proposal for an ISO Threads model, a system consists of multiple threads that execute concurrently within a single Prolog runtime. Threads typically have separate stacks, while large parts of the runtime environment remain shared. While specific dynamic predicates can be declared as thread-local, the default behavior encourages a model where all threads see and modify the same set of predicates and facts. This shared-memory model offers high performance for cooperative tasks but introduces the risk of race conditions and requires the programmer to manage synchronization explicitly using mutexes or message queues. This design gives the programmer substantial freedom, but it also places the burden of correctness on explicit discipline: avoiding races, preventing deadlocks, and ensuring that invariants are preserved under interleavings.

% Web Prolog makes the opposite tradeoff. Actors are intended to be isolated units of computation whose mutable state (if it has one) is encapsulated in the recursive structure of the actor loop or in the actor's private dynamic database, not in shared global facilities. When an actor needs initial programs and/or data, it can be provisioned explicitly at spawn-time. The model therefore encourages a share-nothing architecture by default, reducing the surface area on which shared-memory concurrency bugs can arise.

% Within the taxonomy of programming paradigms defined by Peter van Roy, these two systems represent a fundamental shift in how concurrency and state are handled. ISO Prolog Threads align with the Shared-State Concurrent Programming paradigm. This model combines concurrency with mutable shared state, which van Roy identifies as one of the most complex paradigms to master due to the non-determinism of interleaved access. Web Prolog transitions into the Message-Passing Concurrent Programming paradigm. By eliminating shared mutable state, it moves into a model where interleaving only occurs at the mailbox level. This significantly simplifies formal reasoning and debugging, as each actor's internal logic remains strictly sequential and deterministic.

% \todo{Perhaps move some of this to the van Roy paradigms section?}

% % \subsubsection{Communication and coordination}

% % In the ISO Threads proposal, message passing is one coordination mechanism among several. The draft provides an API for thread creation and termination, message passing via queues, and primitives for mutual exclusion. 

% % In Web Prolog, an actors' mailbox is the central hub for inter-process communication and coordination involving that actor. This is where messages end up when sent by another actor. It is the only queue needed.


% % Sending is asynchronous.

% % Reception is expressed as a multi-clause selective receive. A \texttt{receive/1-2} call can match against several patterns, can skip over non-matching messages without consuming them, and can attach guards expressed as Prolog goals. This allows message acceptance to depend on semantic conditions rather than only structural unification. The result is that prioritization and protocol handling can often be written directly, without auxiliary buffering or dispatch infrastructure.

% % ---

% % The most significant semantic difference appears in message reception. ISO-style \texttt{thread\_get\_message/1} generally operates on a First-In-First-Out (FIFO) basis. While it supports pattern matching, the operation is destructive: once a message unifies, it is consumed. 

% % Web Prolog's \texttt{receive/1-2}, by contrast, implements selective receive. It scans the mailbox for a message matching the provided patterns, skipping unmatched messages while leaving them in the mailbox for later processing. 

% % This allows for complex prioritization logic and guarded reception, where arbitrary Prolog goals can inspect message content before it is officially 'consumed'.


% % \subsubsection{Failure handling}

% % The ISO threads model exposes thread lifecycle primarily through explicit management: developers create threads, arrange for termination, and observe completion by joining or polling, while errors are typically handled as exceptions within the thread that raised them unless additional infrastructure is built to propagate failures across thread boundaries.

% % Web Prolog treats failure as a normal part of concurrent execution and provides coordination mechanisms designed for that reality. Monitoring and linking decouple reliability from send-time checks: instead of attempting to synchronously validate a destination, processes observe termination asynchronously and can react according to an explicit policy. This supports Erlang-style patterns such as supervision, where recovery logic is factored out of the worker and expressed in a separate coordinating process.

% % ---

% % The ISO Threads adopts a strict, defensive approach where communicating with a non-existent thread or queue is often treated as a logic error that must be handled via exceptions. This reflects a system-programming mindset: resources are local and finite, and a phantom destination is a bug. Web Prolog, adopting the conventions of Erlang, uses — send-and-pray — semantics. Sending a message to a non-existent Pid via the \texttt{!/2} operator is a no-op that succeeds silently. In this view, checking for existence before sending is a race condition — the process might terminate immediately after the check — so the system absorbs the uncertainty. Reliability is achieved through monitors, which asynchronously notify the parent if a child terminates.

% % ---

% % A fundamental departure from the ISO model is how Web Prolog handles process failure. In the ISO Threads proposal, error handling is typically synchronous: a parent thread might use \texttt{thread\_join/2} to wait for a child's termination. This approach assumes the parent is actively managing the lifecycle of the child within the same memory space. Web Prolog adopts asynchronous error-handling primitives: monitoring and linking. These mechanisms allow for the construction of supervision hierarchies where failures are propagated as messages, rather than as unhandled exceptions that might crash the entire node.

% % Achieving similar behavior with ISO threads requires significant boilerplate, such as auxiliary watcher threads to perform joins. In Web Prolog, this is reduced to a single primitive: \texttt{monitor(Pid)}. This abstraction is what enables fault-tolerant systems; because the notification is just a message, it works identically whether the failed actor was local or on a server across the globe.

% \subsubsection{Communication and coordination}

% The ISO Prolog Threads draft is intentionally a low-level concurrency substrate. It standardizes predicates for thread management, message queues, and mutual exclusion, and it assumes a shared runtime in which global facilities (notably the Prolog database) are available to all threads. In other words, message passing is supported, but it is only one coordination technique among several, and it lives alongside shared-memory styles that require explicit synchronization.

% Message queues in the ISO model are lightweight and pragmatic. Each thread is associated with a queue addressed by its identifier (or alias), and additional queues can be created when a design calls for explicit channels.  Sending places a copy of a term in a queue and returns immediately.  Receiving is pattern-based: \texttt{thread\_get\_message/1-2} walks the queue and, if necessary, blocks until a term arrives that unifies with the given pattern; a unifying message is then removed, while non-unifying messages remain queued. This already gives a useful form of selectivity, but it remains a single-pattern, single-step receive, and richer protocol logic (priorities, multi-way dispatch, and semantic filtering) must be expressed as user-level control flow around repeated gets, explicit buffering, or auxiliary dispatch threads.

% Web Prolog makes a different commitment. An actor's mailbox is the central communication hub for that actor: all inter-process interaction is expressed as asynchronous sends to a pid and selective reception from that mailbox. In this setting, the mailbox is not merely an available mechanism; it is the primary coordination surface, and it is typically the only queue an application needs to name explicitly. Reception is expressed as a multi-clause \texttt{receive/1-2} whose clauses can be read as a protocol in miniature: several patterns are tried, unmatched messages can be deferred without being consumed, and guards written as Prolog goals can make acceptance depend on semantic conditions rather than structural unification alone. The practical payoff is that prioritization and protocol handling often appear directly in the receive specification, instead of being factored out into ad hoc buffering and dispatch infrastructure.

% \subsubsection{Failure handling}

% The ISO Threads model exposes failure and lifecycle primarily through explicit management and explicit observation. Threads have a well-defined termination status, including the case where the thread ends with an uncaught exception. In the common pattern, a parent creates a thread, allows it to run, and later synchronizes with it (for example, by joining) to discover whether it succeeded, failed, exited, or terminated with an exception. This is consistent with a shared-runtime worldview: the parent is assumed to be able to manage the child's lifecycle within the same address space, and error information is naturally reported at the synchronization points the parent chooses to establish.

% The ISO model is also deliberately strict about invalid endpoints. The draft specifies existence errors when attempting to send to a queue that is not associated with a current queue. This is a reasonable default in a local, resource-oriented model where queues and thread identifiers are treated as concrete runtime objects, and where an invalid destination is typically interpreted as a programming error that should not pass silently.

% Web Prolog treats failure as a routine event in concurrent execution and therefore provides coordination mechanisms that are designed to be used continuously rather than only at join points. In particular, monitoring and linking decouple reliability from send-time validation: rather than synchronously checking whether a destination exists, a process arranges to be notified asynchronously if another process terminates, and it can then react according to an explicit policy. This shifts the developer's attention from defensive endpoint checking (which is inherently racy in a concurrent setting) to explicit failure coordination and recovery structure. It also enables supervision-style designs in which failure handling is factored out of the worker and expressed in a separate coordinating process. Because the notification is itself a message, the same idea scales naturally to remote actors, whereas ISO-style joins are inherently tied to local threads and a shared runtime.

% \subsubsection{Distribution and network transparency}

% The ISO Threads model is, by design, a model of \emph{local} concurrency. Threads share a single Prolog runtime and are intended to coordinate through shared facilities and in-memory synchronization. Distribution can of course be \emph{added}, but only by stepping outside the concurrency abstraction: one must introduce explicit networking libraries, external protocols, and some form of marshalling and routing that is not reflected in the thread primitives themselves. In other words, once computation crosses a machine boundary, the programmer is no longer using ``threads, queues, and mutexes''; they are using a different layer of concepts.

% Web Prolog starts from the opposite end. Its concurrency abstraction is actor-based message passing, and message passing is equally meaningful whether the recipient happens to be local or remote. A pid can denote an actor on the same node or on another node, and the core primitives remain the same: \texttt{spawn/1-3} creates an actor, and \texttt{!/2} sends it a message. Distribution is therefore not a new programming model but a change of \emph{placement} within the same model.

% This is the sense in which Web Prolog aims for \emph{network transparency}: the local and the distributed cases share the same conceptual surface area. The thread model, grounded in shared memory, cannot offer this property as a consequence of its basic assumptions. Put differently, the ISO Threads draft provides the raw materials for building concurrent systems, whereas Web Prolog aims to provide a higher-level engine whose default mode of operation already matches the distributed setting of the Web.


% % \subsubsection{Distribution and network transparency}

% % Perhaps the most consequential difference is that the ISO threads model is inherently local. Threads share a runtime and rely on shared-memory assumptions, so distribution must be introduced separately via libraries and protocols that live outside the concurrency abstraction.

% % Web Prolog's actor model is designed to scale from local to remote execution without changing the programming model. The same \texttt{spawn} and send primitives remain meaningful when the callee is remote. In other words, distribution does not require a new concurrency model — it is a change of placement within the same model.

% % \todo{Place the ISO Threads model and the actor model within van Roy's distinction between programming in the small and programming in the large.}

% % ---

% % Furthermore, Web Prolog achieves network transparency. The same \texttt{spawn} and \texttt{send} primitives used for local processes work across nodes — a feat impossible with a shared-memory thread model. While the ISO standard provides the raw materials for concurrency, Web Prolog provides a high-level engine designed for the distributed web.


% \subsubsection{ISO Threads as substrate}

% Several Prolog systems (e.g. SWI-Prolog, YAP, and XSB) implement substantial parts of the ISO Threads proposal, and the concrete predicate set varies somewhat across implementations. 

% It is important to emphasize that the ISO Threads standard and the Web Prolog Actor Model are not mutually exclusive; rather, they operate at different levels of the system stack. In many implementations, the ISO \texttt{thread\_*} predicates serve as the very foundation upon which a Web Prolog node is built. While the actor model provides the high-level paradigm for application logic and distribution, the underlying engine uses the ISO thread model to manage the physical execution of those actors on a multi-core machine.

% This relationship reveals the strategic importance of the shared global database in the ISO model. In a Web Prolog node, the shared database acts as the primary repository for the node's core knowledge base, typically provided and managed by the node's owner. While Web Prolog enforces isolation to prevent client actors from interfering with one another's internal state, these actors often require access to the node's collective intelligence. In this architecture, client actors treat the global database as a read-only resource. This allows many concurrent actors to perform high-performance queries against a massive, shared set of facts and rules — such as a large-scale ontology or a dataset — without the memory overhead of duplicating that data for every spawned process.

% By leveraging ISO Threads as the implementation layer, Web Prolog gains the best of both worlds: the raw efficiency of shared-memory access for common knowledge, and the 'saner' safety of the actor model for concurrent coordination and state management. The ISO standard gives the system architect the tools to build the engine, while the actor model gives the application developer a robust, network-transparent environment in which to run it. 


% \subsubsection{Summary}

% The ISO threads draft provides a flexible low-level substrate: threads, mutexes, and queue-based communication. Web Prolog deliberately restricts the space of concurrency mechanisms in order to make a particular paradigm — isolated processes communicating by message passing, with explicit failure coordination — the default. The argument of this book is that this restriction buys conceptual clarity, compositional reliability patterns, and a smooth path from local concurrency to distributed programming on the Web.

% The practical implications of these design choices are most visible in code volume. High-level concurrency predicates like \texttt{concurrent/2-3} in Web Prolog often require significantly less code than an equivalent ISO Threads implementation. This reduction in boilerplate is not merely aesthetic; it reduces the cognitive load on the developer and the surface area for bugs. Smaller codebases enhance the clarity of intentions, which is critical for long-term projects.
